% ============================================================
% Scaling Laws, Class Invariance, and the Limits of Audibility
% of an Explicit Prime-Zero Coupling Operator
%
% Author:   Ulrich Tehrani
% Date:     August 2026
% Version:  v0.30
% Zenodo:   10.5281/zenodo.21899170
% License:  CC BY 4.0
% ============================================================

\documentclass[11pt]{article}
\usepackage{cmap}
\usepackage[T1]{fontenc}
\usepackage[utf8]{inputenc}
\usepackage[a4paper,margin=1in]{geometry}
\usepackage{amsmath,amssymb,amsthm,mathtools}
\usepackage{xcolor}
\usepackage{hyperref}
\usepackage{enumitem}
\usepackage{setspace}
\usepackage{booktabs}
\usepackage{array}
\usepackage{graphicx}

\hypersetup{colorlinks=true,
  linkcolor=blue!50!black,
  urlcolor=blue!50!black,
  citecolor=blue!50!black,
  pdftitle={Scaling Laws, Class Invariance, and the Limits of Audibility
            of an Explicit Prime-Zero Coupling Operator},
  pdfauthor={Ulrich Tehrani}}
\setlength{\parskip}{0.5em}
\setlength{\parindent}{0pt}
\onehalfspacing
\emergencystretch=1em

% --------------------------------------------------------
% Theorem environments (numbered within section)
% --------------------------------------------------------
\newtheorem{theorem}{Theorem}[section]
\newtheorem{proposition}[theorem]{Proposition}
\newtheorem{lemma}[theorem]{Lemma}
\newtheorem{corollary}[theorem]{Corollary}
\newtheorem{definition}[theorem]{Definition}
\newtheorem{assumption}[theorem]{Assumption}
\newtheorem{observation}[theorem]{Observation}
\newtheorem{conjecture}[theorem]{Conjecture}
\newtheorem*{remark}{Remark}
\newtheorem{openproblem}[theorem]{Open Problem}

% --------------------------------------------------------
% Series-wide macros
% --------------------------------------------------------
\newcommand{\Hstr}{H_{\mathrm{str}}}
\newcommand{\Hnull}{H_{\mathrm{null}}}
\newcommand{\Tt}{\widetilde{T}}
\newcommand{\DSEL}{D_{\mathrm{SEL}}}

% --------------------------------------------------------
% Paper-specific macros
% --------------------------------------------------------
\newcommand{\Dclass}{\mathcal{D}_\delta}        % Beurling density class
\newcommand{\dBL}{d_{\mathrm{BL}}}              % bounded-Lipschitz distance
\newcommand{\Ddisc}{\Delta_{\mathrm{disc}}}     % certified discrimination margin
\newcommand{\WeN}{W_{\varepsilon,N}}
\newcommand{\Bsm}{B_{\mathrm{sm}}}              % curvature of the smooth control world

% --------------------------------------------------------
% Title
% --------------------------------------------------------
\title{\vspace{-1.5em}\bfseries Scaling Laws, Class Invariance, and the
Limits of Audibility\\[0.4em]
\large of an Explicit Prime--Zero Coupling Operator}
\author{Ulrich Tehrani}
\date{Research Note \textperiodcentered\ August~2026
      \textperiodcentered\ v0.30}

% ============================================================
\begin{document}
\maketitle
% ============================================================

\begin{abstract}
This paper determines what a prime-indexed coupling operator can and cannot
hear. The operator $\Tt(\sigma)=\Phi(\sigma)\Phi(\sigma)^*$ of the present
series is built explicitly from prime phasors; no spectral property is
postulated. Three questions are settled about it, and the order in which they
are settled matters. A symmetry barrier comes first: every symmetric
functional of the orbit of a zero under the functional equation is even in the
signed distance from the critical line, so the class of observables studied
here is blind to the left/right orientation of that displacement. Evenness
needs no regularity; under differentiability the first-order response
vanishes, and for the curvature sums of this series, whose dependence enters
through even orbit averages of the entire weighted curvature test kernel, the
first
possible non-constant term is quadratic; it is the first surviving term unless the
aggregate quadratic coefficient cancels. This is a theorem, not a limitation
of effort. What
remains is the audible part --- audibility always meaning separation from
the stated comparison worlds by the stated observable class --- and it splits
cleanly. On the prime side, what
is audible is stated exactly as certified: the prime-side weighted point
configuration is compared with three control objects kept separate
throughout --- a smooth measure whose own zeta function is zero-free, its
rescaling to the exact mass of the prime counting measure, and a discrete
equal-mass quadrature world. Certified separation margins are established
against the smooth measure and the discrete world; only the rescaling
matches the prime count exactly;
and for the discrete pair no functional of the counting data with
bounded-Lipschitz constant below an explicit threshold factors the
separation. This identifies the point configuration relative to the stated
controls, not multiplicativity as its cause: the controls vary point geometry
and multiplicative structure together. On the prime input side, the scaling laws are invariant across the Beurling
class: every admissible system carries the same limiting function, the same
certified separation and the same leading profile with the same numerical
band --- a within-class non-distinction, proved, not a separation. On the
ordinate side, a pre-registered unfolded discrimination test over five worlds
and five observables finds no separation beyond the imposed unit-density
skeleton, and the scope of that measurement is declared.

\textbf{What is proved.}
The kernel identity representing the companion prime-space Gram operator
$T=\Phi^{*}\Phi$ as the damped inverse spectrum evaluated at prime ratios
and products (Theorem~\ref{thm:kernel});
canonicity of the phasor kernel within the unimodular trigonometric almost
periodic class (Lemma~\ref{lem:canonicity},
Proposition~\ref{prop:resonance});
non-convergence of the normalised leading spectral ratio at the reference
data $(100,0.05)$, with a certified two-point separation
(Theorem~\ref{thm:nonconv});
divergence of the mean energy-asymmetry functional along a bounded leading profile
(Theorem~\ref{thm:eta});
the ordinate tail bound, which shows that the ordinate count is a truncation
and not a limit parameter (Theorem~\ref{thm:sat});
transfer of both scaling laws to Beurling systems, carrying the same
limiting objects and the same leading profile class-independently
(Theorem~\ref{thm:class}) together with the sideband transfer
(Theorem~\ref{thm:sideband});
the summation formula for the smooth measure (Theorem~\ref{thm:inv}), the
quadrature estimate for the discrete world (Theorem~\ref{thm:quad}) and the
factorisation barrier (Theorem~\ref{thm:fact});
and the symmetry barrier (Theorem~\ref{thm:barrier}).

\textbf{What is certified.}
The two-point separation of the limiting ratio; the discrimination margins in
curvature against the smooth control world, uniformly on the reference window;
and the bounded-Lipschitz distance between the two weighted counting measures.
Each certified statement names its verification protocol where it is first
stated.

\textbf{What is numerical.}
The outcome of the pre-registered unfolded discrimination test
(\S\,\ref{sec:audibility}), reported as a measurement with its design, its
criterion, its result and its scope; the completeness counts underlying it;
the embedding of the sideband transfer in an explicit test world; and the
robustness of the second moment under completion of the mode set to prime
powers (Remark~\ref{rem:lambda}); and a function-side reconstruction
experiment, separate from the battery, reported with its design, its
metrics and its scope.

\textbf{What remains open.}
Pointwise, as opposed to mean, finality of the divergence
(\S\,\ref{sec:open}); the minimal order break that destroys the class laws;
the exceptional-set measure under intermediate conditions; extension of the
factorisation barrier from one pair of worlds to families; the
identification design that would attribute the certified prime-side
separation to a single feature; pair and correlation statistics on the
unfolded ordinate axis; and the combinatorial gap that separates the
canonicity lemma from the full Bohr class.

\textbf{Structure.}
The formal core is \S\S\,2--6, and it is non-circular: every statement there
carries its status in its heading --- proved, certified, or proved from
certified inputs --- and numerical observations are labelled as such where
they occur. No use of the Riemann Hypothesis, the GUE conjecture, or the
Hilbert--P\'olya postulate is made anywhere in \S\S\,2--6.
\end{abstract}

% --------------------------------------------------------
% Notation and Conventions
% --------------------------------------------------------
\section*{Notation and Conventions}
\label{sec:notation}

\emph{Critical line.}
Throughout, $\sigma$ denotes $\Re(s)$ and the critical line means
$\sigma=\tfrac12$. The trace parameter of the operator family is likewise
written $\sigma$ and is evaluated at $\sigma_0=\tfrac12$ unless stated
otherwise.

\medskip
\emph{Global parameters.}
The reference values are $\kappa=53$ (prime cutoff, giving $16$ active primes
$p\le\kappa$), $\varepsilon=0.05$ (Gaussian regularisation parameter), and
$N=100$ (number of Riemann ordinates $\gamma_k$ employed). Certified
statements are given on the window $\varepsilon\in[0.04,0.07]$. The unfolded
scale used in \S\,\ref{sec:audibility} is $\tilde\varepsilon=0.25$; it is
derived, not chosen, from the effective bandwidth at the reference parameters.

\medskip
\emph{Main objects.}
The coupling map $\Phi(\sigma):\Hstr\to\Hnull$ between the finite-dimensional
prime space $\Hstr=\ell^2(\{p\le\kappa\})$ and the ordinate space
$\Hnull=\ell^2(\{\gamma_1,\dots,\gamma_N\})$ is
\[
\Phi(\sigma)_{k,p}=e^{-\varepsilon^2\gamma_k^2/2}\,\sin(\sigma\gamma_k\log p),
\]
and the coupling operator is $\Tt(\sigma)=\Phi(\sigma)\circ\Phi(\sigma)^*$ on
$\Hnull$, with loop operator $T(\sigma)=\Phi(\sigma)^*\circ\Phi(\sigma)$ on
$\Hstr$. The oscillatory trace sum and the curvature sum are
\[
O(\sigma)=\tfrac12\sum_{k,p}e^{-\varepsilon^2\gamma_k^2}
          \cos(2\sigma\gamma_k\log p),
\qquad
B=\sum_{k,p}e^{-\varepsilon^2\gamma_k^2}(\gamma_k\log p)^2
  \cos(\gamma_k\log p),
\]
the second being the on-line curvature sum with $O''(\tfrac12)=-2B$. The
damped inverse spectrum entering \S\,\ref{sec:operator} is
\[
G(u)=\sum_{k=1}^{N}e^{-\varepsilon^2\gamma_k^2}\cos(\sigma\gamma_k u),
\]
a single even function of one real variable. The Beurling density class of
\S\,\ref{sec:class} is
\[
\Dclass=\{\mathfrak{B}:\ \vartheta_{\mathfrak{B}}(x)=x+O(x/(\log x)^{2+\delta})\},
\qquad\delta>0.
\]

\medskip
\emph{Weights.}
Two prime weights occur in this series and are numerically distinct. With
\[
V_p(\tfrac12)=\frac{4(\log p)^2p}{(p-1)^2},
\qquad
f_p=V_p(\tfrac12)-\frac{4(\log p)^2}{p}=\frac{4(\log p)^2(2p-1)}{p(p-1)^2},
\]
the
$\eta$-framework weight is $c_p^{\eta}:=\sqrt{V_p(\tfrac12)}$ and the
renormalised Weil weight is $c_p^{\mathrm{ren}}:=\sqrt{f_p}$. At the reference
parameters the associated values of the energy-asymmetry functional are
$\eta_{\mathrm{orig}}(\tfrac12)=0.69078176$ and
$\eta_{\mathrm{ren}}(\tfrac12)=0.66926873$. In the convention of the earlier
papers $\eta[c]=1-E_{\mathrm{spec}}/E_{\mathrm{str}}$, so that
$E_{\mathrm{spec}}/E_{\mathrm{str}}=1-\eta$ is the Rayleigh quotient of that
series; $\eta\to-\infty$ therefore means
$E_{\mathrm{spec}}/E_{\mathrm{str}}\to+\infty$. Both are evaluated in the
convention in which the phasor argument carries no trace parameter, that is
$\sin(\gamma_k\log p)$; this is the convention of the earlier papers of the
series in which these quantities were introduced, and it does not coincide
with $\sin(\sigma\gamma_k\log p)$ at $\sigma=\tfrac12$.

\medskip
\emph{Analytical foundations.}
The proofs use the prime number theorem with classical error term together
with Mertens' theorem, Weyl's eigenvalue perturbation inequality, Dirichlet's
pigeonhole principle, and, at exactly one place, the two-function Parseval
identity for Besicovitch $B^2$ almost periodic functions. The numerical
completeness checks of \S\,\ref{sec:audibility} use the argument principle;
it enters no proof.

\medskip
\emph{Numerical computations.}
Working precisions vary by artefact and are declared here: the five-world
battery runs at $40$ digits with \texttt{float64} observable paths; the
function-side reconstruction at $8$ digits with NumPy/SciPy; the two margin
and separation certificates in Arb ball arithmetic at $192$ and $320$ bits
through \texttt{python-flint}, with an \texttt{mpmath} cross-layer at $32$
digits. Riemann ordinates $\gamma_k$ enter through the normative layer of
\S\,\ref{sec:methods}, regenerated via \texttt{mpmath.zetazero} or
isolated in Arb, in either case under exact double-precision cross-check. Verification code is
available at \url{https://github.com/utehrani/analysislab-nt}.

\medskip
\emph{Not to be confused with.}
The following pairs are distinct objects and are never interchanged.

\begin{center}
\begin{tabular}{lll}
\toprule
Symbol & Meaning & Distinct from \\
\midrule
$G(u)$ & damped inverse spectrum & $g$, the candidate kernel of
  \S\,\ref{sec:operator} \\
$\mathcal{F}$ & Bohr spectrum of $g$ & $\Lambda(n)$, von Mangoldt \\
$\omega_0$ & unimodular constant phase & $\zeta(s)$ \\
$\mathcal{O}$ & observable (functional on worlds) & $O(\sigma)$, trace sum \\
$\DSEL\approx 10.9845$ & trace-formula constant & $D\approx 9.470$ \\
$c_p^{\eta}$ & $\eta$-framework weight & $c_p^{\mathrm{ren}}$, Weil weight \\
$B$ & curvature sum, $\gamma_k^2$ weights &
  $B_{\mathrm{int}}^{+}(\varepsilon,N)$, integrated bias \\
$r(x)$ & normalised leading spectral ratio & bias ratios $r_p$ of the earlier
  papers \\
$\mathcal K_{\kappa,\varepsilon}(z)$ & aggregated curvature test kernel &
  $G(u)$ (ordinate sum); $K_\varepsilon$ (\S\,5) \\
$\varphi$, $K_\varepsilon$ & test function, curvature kernel & $g$, the
  candidate kernel of \S\,\ref{sec:operator} \\
$\mathfrak{B}$ (\S\,4) & Beurling system index & $B$, $\Bsm$, $B[A]$, $B_{\mathrm{int}}^{+}$; world $A$ (\S\,5) \\
\bottomrule
\end{tabular}
\end{center}

The curvature sums attached to different point sets are distinguished
throughout by their argument: $B[P]$ for the primes $p\le\kappa$, $B[A]$ for
the discrete quadrature world $A$, and $\Bsm$ for the smooth measure
$\mu_{\mathrm{sm}}$ of \S\,\ref{sec:barrier}. The three control objects of
that section --- $\mu_{\mathrm{sm}}$, its mass-matched rescaling
$\mu_{\mathrm{norm}}$, and the discrete world $A$ --- are kept apart
throughout and are never referred to collectively as ``the control world''.

\vspace{2em}
\newpage

% ============================================================
% §1 Introduction
% ============================================================
\section{Introduction}
\label{sec:intro}

% --------------------------------------------------------
\subsection{Background and Motivation}
% --------------------------------------------------------

This paper is the ninth in a series developing an operator-theoretic framework
for studying the distribution of the non-trivial zeros of the Riemann zeta
function. It closes the series.

\textbf{Paper~1}~\cite{Paper1} decomposed the explicit formula into a
curvature structure and isolated the prime-local contributions that carry it.

\textbf{Paper~2}~\cite{Paper2} constructed the renormalised weights
$c_p^{\mathrm{ren}}$ and connected the local curvature to the Weil functional.

\textbf{Paper~3}~\cite{Paper3} introduced the finite-cutoff Hilbert-space
model, the prime space $\Hstr$, the ordinate space $\Hnull$, and the
energy-asymmetry functional $\eta_{\mathrm{orig}}$ measuring the energy defect
between them.

\textbf{Paper~4}~\cite{Paper4} exhibited the dual operator and found, on
its tested finite-cutoff grid, no numerical support for a Hilbert--P\'olya
interpretation: the tested eigenvalues are not the ordinates, and the
inverse-spectrum correlation degrades with cutoff. No structural
obstruction theorem was proved there.

\textbf{Paper~5}~\cite{Paper5} proved the spectral trace formula
$\mathrm{Tr}\,\Tt(\sigma)=\DSEL-O(\sigma)$ and fixed the constant $\DSEL$.

\textbf{Paper~6}~\cite{Paper6} established the curvature identity
$O''(\tfrac12)=-2B$ at the critical line.

\textbf{Paper~7}~\cite{Paper7} analysed stationarity of $O'$ under an
explicitly named logarithmic-saving hypothesis, which asks only for a
logarithmic saving in the ordinate aspect relative to a trivial bound and is
not known to follow from the Riemann Hypothesis.

\textbf{Paper~8}~\cite{Paper8} proved unconditionally that the second-moment
bias of the full Guinand--Weil object is negative, and controlled hypothetical
off-critical zeros using the verified height of Platt and
Trudgian~\cite{PlattTrudgian21}.

\textbf{The present paper} asks the question that characterises the object
rather than extending it: what does this coupling hear, and what can it
provably not hear? The answer is a classification along three axes, and the
barrier axis is stated first because it governs how the other two may be read.

The operator-theoretic approach to the zeros has a long history, and the
present construction sits at a definite distance from its main lines. In the
programme of Connes and collaborators the spectral realisation is adelic and
the primes enter through truncated Euler products, that is through a prime
\emph{cutoff}; recent work in that direction~\cite{Connes26,ConnesVS25} and
its independent numerical implementations~\cite{Groskin26a} keep this
architecture. Here the primes are instead the \emph{index set} of the
operator, and the ordinates enter as frequencies; the difference is
architectural and is set out in \S\,\ref{sec:discussion}. The
semiclassical picture of Berry and Keating, and the random-matrix statistics
descending from Montgomery~\cite{Montgomery73} and
Odlyzko~\cite{Odlyzko87} through
Bogomolny and Keating~\cite{BogomolnyKeating96},
Rudnick and Sarnak~\cite{RudnickSarnak96} and
Katz and Sarnak~\cite{KatzSarnak99}, provides the background motivation for
comparing unfolded worlds in \S\,\ref{sec:audibility}; it is neither an input
nor a hypothesis tested by the registered battery.

% --------------------------------------------------------
\subsection{Main Results}
% --------------------------------------------------------

\S\,\ref{sec:operator} records the construction and the two identities
imported from the earlier papers, and adds two results published here for the
first time. The kernel identity
$T_{p,q}=\tfrac12[G(\log(p/q))-G(\log(pq))]$ shows that the companion
prime-space Gram operator $T=\Phi^{*}\Phi$ \emph{is} the damped inverse
spectrum sampled at prime ratios and prime products: one function of one
variable, contracted in two ways. The canonicity
lemma establishes character rigidity: within the unimodular trigonometric
almost periodic class with pairwise orthogonal dilates, the building block is
forced to be a character. The frequency assignment $c=\sigma\gamma_k$ and the
passage to the odd part $\sin$ are recorded conventions
(\S\,\ref{sec:operator}).

\S\,\ref{sec:scaling} proves that, at the reference data $(100,0.05)$, the
normalised leading spectral ratio does not converge as the cutoff grows,
with a certified two-point separation
$\Delta_0\ge 0.0840112$; that the energy-asymmetry functional
$\eta_{\mathrm{orig}}(\kappa)$ diverges to $-\infty$ in mean, an exhaustive
scan over prime cutoffs $\kappa\le2\cdot10^{5}$ finding the last positive
value at $\kappa=146521$ and a negative one at the next prime cutoff
$146527$; and that the ordinate count
enters only through an explicit tail bound, so that the genuine limit
parameters are $(\kappa,\varepsilon)$ and not $(\kappa,N)$.

\S\,\ref{sec:class} shows that both scaling laws are invariant across the
Beurling class $\Dclass$: they hold over every admissible system, with the
\emph{same} limiting function, the \emph{same} certified separation and the
\emph{same} leading profile, whose numerical band is class-independent --- a
within-class non-distinction, not an ordinate-side separation. What these
observables provably retain of the prime input, at the level of the proved
leading laws, is the density class: $R$, $\Delta_0$, $F$ and the band are
class-independent. Lower-order, system-dependent contributions are not
excluded. A sideband transfer
makes the response precise for a different object, a smooth generalised-prime
measure with a bounded oscillating density: the laws are not destroyed, the
limiting objects move by explicit replacements.

\S\,\ref{sec:barrier} contrasts the primes with three control objects, kept
separate throughout: the smooth measure $\mu_{\mathrm{sm}}$, whose own zeta
function is zero-free; its rescaling $\mu_{\mathrm{norm}}$, which is the one
matching the prime count exactly; and the discrete equal-mass quadrature world
$A$. A summation formula and a quadrature estimate control the main terms; the
curvature does not follow, and the discrepancy against $\mu_{\mathrm{sm}}$ and
against $A$ is certified uniformly on the reference window. No functional of
the counting data with bounded-Lipschitz constant below $35.51$ can factor the
discrepancy against $A$.

\S\,\ref{sec:audibility} states the symmetry barrier and reports the
pre-registered unfolded discrimination test as a measurement.
\S\,\ref{sec:discussion} assembles the classification and delineates it
against neighbouring programmes. \S\,\ref{sec:methods} documents the ordinate
protocol and the certificate stack; \S\,\ref{sec:open} leaves the open
questions.

% --------------------------------------------------------
\subsection{Reader Contract}
% --------------------------------------------------------

This paper \emph{proves} the kernel identity and the canonicity of the phasor
kernel (\S\,\ref{sec:operator}), the scaling laws, stated at the
reference data $(100,0.05)$ (\S\,\ref{sec:scaling}), the class transfer with the same limiting objects and the
sideband transfer for a smooth oscillating measure (\S\,\ref{sec:class}), the main-term equality and the
factorisation barrier --- the latter proved from certified inputs
(\S\,\ref{sec:barrier}) --- and the symmetry barrier
(\S\,\ref{sec:audibility}). These are unconditional mathematical statements.

This paper \emph{establishes numerically} the outcome of the pre-registered
discrimination test, the completeness counts supporting it, the embedding of
the sideband transfer in an explicit test world, the robustness of the second
moment under completion of the mode set to prime powers, and the
function-side reconstruction experiment of \S\,\ref{sec:discussion}. These
statements are verified on finite grids and are labelled \emph{numerical} in
their headings.
A smaller group of statements is labelled \emph{certified}: they are bounds
obtained by interval arithmetic with an explicit error protocol, named where
they occur.

This paper \emph{leaves open} the questions collected in \S\,\ref{sec:open}:
pointwise finality of the divergence, the minimal order break, the
exceptional measure under intermediate conditions, the family form of the
factorisation barrier, the identification design for the prime-side
separation, pair statistics on the unfolded axis, and the combinatorial gap
in the canonicity lemma.

Results imported from the earlier papers are marked \emph{proved (cited)} and
are not reproved here; the citation names the paper and the section in which
the proof stands.

% ============================================================
\section{The Coupling Operator: Construction, Trace Identity, and Kernel}
\label{sec:operator}
% ============================================================

The construction is explicit and finite. Fix $\kappa$, $\varepsilon$, $N$. The
matrix $\Phi(\sigma)$ has one column per prime $p\le\kappa$ and one row per
ordinate $\gamma_k$, and $\Tt(\sigma)=\Phi\Phi^*$, $T(\sigma)=\Phi^*\Phi$ are
the two Gram operators attached to it. Since $\Phi$ is a finite real matrix,
$T$ and $\Tt$ are positive semidefinite and

\[
\mathrm{spec}(T)\setminus\{0\}=\mathrm{spec}(\Tt)\setminus\{0\};
\]
 nothing is
postulated about either spectrum.

\begin{theorem}[Trace formula and curvature identity; proved (cited)]
\label{thm:imported}
At the reference parameters,
$\mathrm{Tr}\,\Tt(\sigma)=\DSEL-O(\sigma)$ with $\DSEL=10.9845$
(\cite{Paper5}, \S\,3), and $O''(\tfrac12)=-2B$ (\cite{Paper6}, \S\,6).
Numerically $B=-19342.5476$, hence $O''(\tfrac12)=+38685.0952$, while the
first derivative is small but not zero, $O'(\tfrac12)=+2.4751$, in the
three-term decomposition of that derivative (\cite{Paper7}, \S\,4); the ratio
of the two quantifies proximity to stationarity without establishing it. The second-moment curvature of the full
Guinand--Weil object is negative unconditionally at the reference cutoff
(\cite{Paper8}, \S\S\,5--6).
\end{theorem}

The first new observation is that the prime-space Gram operator
$T=\Phi^{*}\Phi$ is an arithmetic kernel evaluated at prime ratios and
products. The ordinate-space companion $\Tt=\Phi\Phi^{*}$ has the same
non-zero spectrum, but the following identity is an entrywise formula for
$T$, not for $\Tt$.

\begin{theorem}[Kernel identity; proved]
\label{thm:kernel}
For all primes $p,q\le\kappa$ and all $\sigma$,
\[
T_{p,q}
=\tfrac12\bigl[\,G(\log(p/q))-G(\log(pq))\,\bigr],
\qquad
G(u)=\sum_{k=1}^{N}e^{-\varepsilon^2\gamma_k^2}\cos(\sigma\gamma_k u).
\]
The identity is algebraic and exact; it holds entrywise, for every choice of
ordinate set and every regularisation width.
\end{theorem}

\begin{proof}
By definition

\[
T_{p,q}=\sum_k e^{-\varepsilon^2\gamma_k^2}
\sin(\sigma\gamma_k\log p)\sin(\sigma\gamma_k\log q).
\]

The product formula
$\sin a\sin b=\tfrac12[\cos(a-b)-\cos(a+b)]$ with
$a=\sigma\gamma_k\log p$, $b=\sigma\gamma_k\log q$ gives
$a-b=\sigma\gamma_k\log(p/q)$ and $a+b=\sigma\gamma_k\log(pq)$. Summing over
$k$ and using the definition of $G$ yields the statement.
\end{proof}

Thus the single function $G$ determines every entry of the prime-space Gram
matrix through its values at $\log(p/q)$ and $\log(pq)$: the first channel
measures logarithmic differences, the second logarithmic sums. One consequence
is a resolution scale. Two primes whose logarithms are close are seen by $G$
through nearly the same argument, and their columns become nearly parallel.
\emph{Numerical observation} (resolution diagnostics). At
the reference parameters the Gram correlations of the twin pairs $(41,43)$ and
$(29,31)$ reach $0.937$ and $0.930$ respectively, while distant pairs such as
$(47,53)$ and $(2,3)$ give $0.134$ and $-0.521$. Counting pairs above the $0.9$ correlation
threshold --- the criterion used for fusion throughout --- against the
sixteen active primes leaves fourteen resolvable prime directions at the
reference parameters.
The correlation is not monotone in the logarithmic distance, because the
product term $G(\log(pq))$ modulates it.

\subsection{Canonicity of the phasor kernel}

The kernel $\sin(\sigma\gamma_k\log p)$ was introduced in the earlier papers
as a construction. What is forced --- and all that is forced --- is its
character part; the remaining steps are conventions, and they are recorded as
such below. Write $M[f]$ for the Bohr mean of an almost periodic $f$.

\begin{definition}[Bohr mean and dilate orthogonality]
For a Bohr almost periodic function $f:\mathbb{R}\to\mathbb{C}$ write

\[
M[f] := \lim_{T\to\infty} \tfrac1T \int_0^T f(x)\,dx
\]
 for its Bohr mean,
and 
\[
g\sim \sum_{\lambda\in\mathcal{F}} a_\lambda e^{i\lambda x}
\]
 for its
Bohr--Fourier series, where $\mathcal{F} = \{\lambda : a_\lambda \neq 0\}$ is
the (countable) spectrum and $\sum_{\lambda}|a_\lambda|^2 = M[|g|^2]$
(Parseval). We call the dilate family $\{g(\alpha\,\cdot)\}_{\alpha>0}$
\emph{pairwise orthogonal} if

\[
M_x\!\left[g(\alpha x)\,\overline{g(\beta x)}\right] = 0
\]

for all $\alpha \neq \beta > 0$.
\end{definition}

Here $g$ is a candidate kernel, a unimodular almost periodic function --- not
the damped kernel $G$ of Theorem~\ref{thm:kernel}, which is neither unimodular
nor a candidate.

\begin{lemma}[Character rigidity under unimodularity and dilate
orthogonality; proved]
\label{lem:canonicity}
Let 
\[
g(x) = \sum_{\lambda\in\mathcal{F}} a_\lambda e^{i\lambda x}
\]
 be a
trigonometric polynomial ($\mathcal{F}\subset\mathbb{R}$ finite, all
$a_\lambda\neq 0$) with
\begin{enumerate}
\item[(H1)] $|g(x)| = 1$ for all $x\in\mathbb{R}$ (pointwise unimodularity),
\item[(H2)] 
\[
M_x\!\left[g(\alpha x)\,\overline{g(\beta x)}\right] = 0
\]

  for all $\alpha\neq\beta>0$ (dilate orthogonality).
\end{enumerate}
Then $g(x) = \omega_0\, e^{icx}$ with $|\omega_0| = 1$ and a uniquely determined
frequency $c \in \mathbb{R}\setminus\{0\}$. Conversely, every such
character satisfies (H1)--(H2); in particular $e^{icx}$ and $e^{-icx}$
both do, so the hypotheses determine the character only as a
$\pm c$-pair, while $c$ is unique for each individual $g$.
\end{lemma}

\begin{proof}
\emph{Step 1 (unimodularity kills multi-frequency spectra).}
Expanding the finite sum,

\[
|g(x)|^2 = \sum_{\lambda,\mu\in\mathcal{F}} a_\lambda \overline{a_\mu}\,
e^{i(\lambda-\mu)x} \equiv 1.
\]

By uniqueness of the coefficients of a trigonometric polynomial, for every
$\tau\neq 0$,
\begin{equation}\label{eq:addAC}
\sum_{\substack{\lambda,\mu\in\mathcal{F}\\ \lambda-\mu=\tau}}
a_\lambda \overline{a_\mu} \;=\; 0 .
\end{equation}
Suppose $\#\mathcal{F} \geq 2$ and set $\lambda^{*} := \max\mathcal{F}$,
$\lambda_{*} := \min\mathcal{F}$, $\tau_0 := \lambda^{*}-\lambda_{*} > 0$.
The only pair $(\lambda,\mu)\in\mathcal{F}^2$ with $\lambda-\mu=\tau_0$ is
$(\lambda^{*},\lambda_{*})$, since $\lambda \le \lambda^{*}$ and
$\mu \ge \lambda_{*}$ force equality in both. Hence \eqref{eq:addAC} at
$\tau=\tau_0$ reads $a_{\lambda^{*}}\overline{a_{\lambda_{*}}} = 0$,
contradicting $a_\lambda\neq 0$ on $\mathcal{F}$. Thus $\mathcal{F} = \{c\}$ and,
by Parseval, $|a_c|^2 = M[|g|^2] = 1$, i.e.\ $g = \omega_0 e^{icx}$ with
$|\omega_0|=1$.

\emph{Step 2 (orthogonality excludes the constant).}
If $c = 0$ then $g \equiv \omega_0$ and

\[
M[g(\alpha x)\overline{g(\beta x)}] = |\omega_0|^2 = 1 \neq 0
\]
 for every
$\alpha\neq\beta$, contradicting (H2). Hence $c\neq 0$.

\emph{Converse.} For $g = \omega_0 e^{icx}$ with $c\neq0$ and
$\alpha\neq\beta>0$ one has
$M[e^{ic(\alpha-\beta)x}] = 0$ because $c(\alpha-\beta)\neq 0$, and
$|g|\equiv 1$ trivially.
\end{proof}

\begin{remark}[Division of labour between the hypotheses]
Neither hypothesis is redundant, and within the trigonometric class their
roles are sharply asymmetric. $g=\cos(cx)$ satisfies (H2) --- the dilate
resonances $c\alpha=\pm c\beta$ require $\alpha=\beta$ or $\alpha=-\beta$,
both excluded --- but violates (H1) and is not a character: (H2) alone is
insufficient. Conversely $g(x)=e^{i\sin x}$ satisfies (H1) with infinite
spectrum $\mathbb{Z}$ (Jacobi--Anger, coefficients $J_n(1)$; identity
$\sum_n J_n(1)J_{n+k}(1)=\delta_{k0}$ realises \eqref{eq:addAC}), yet
violates (H2): for irrational $\beta/\alpha$ the only dilate resonance is
$(0,0)$, giving 
\[
M[g(\alpha x)\overline{g(\beta x)}] = J_0(1)^2 \approx
0.5855 \neq 0.
\]
 So (H1) alone is insufficient once the spectrum is
infinite. Inside the trigonometric class, Step~1 shows that (H1) alone
already forces a single frequency; (H2) contributes exactly the exclusion
of the constant character.
\end{remark}

\begin{proposition}[Resonance identities in the full Bohr class; proved]
\label{prop:resonance}
Let $g$ be Bohr almost periodic with Bohr--Fourier series

\[
\sum_{\lambda\in\mathcal{F}} a_\lambda e^{i\lambda x}.
\]
 Then for all
$\alpha,\beta>0$,
\begin{equation}\label{eq:multAC}
M_x\!\left[g(\alpha x)\,\overline{g(\beta x)}\right]
= \sum_{\substack{\lambda,\mu\in\mathcal{F}\\ \alpha\lambda=\beta\mu}}
a_\lambda\overline{a_\mu}.
\end{equation}
If in addition $|g|\equiv 1$ pointwise (H1), then \eqref{eq:addAC} holds
for every $\tau\neq 0$, and every difference $\tau\neq0$ realised in
$\mathcal{F}$ is realised by at least two distinct pairs. If, moreover, the
dilates of $g$ are pairwise orthogonal (H2), then:
\begin{enumerate}
\item[(a)] $0\notin\mathcal{F}$: by (H2) the left-hand side of
\eqref{eq:multAC} vanishes for all $\alpha\neq\beta$; choosing
$\beta/\alpha$ outside the countable set

\[
\{\lambda/\mu : \lambda,\mu\in\mathcal{F},\ \mu\neq0\}\cup\{1\}
\]
 leaves
$(0,0)$ as the only resonance on the right-hand side, forcing $|a_0|^2=0$.
\item[(b)] Ratio rigidity: every ratio $r\neq1$ realised by a same-sign
pair in $\mathcal{F}$ is realised by at least two distinct pairs (otherwise
\eqref{eq:multAC} with $\beta/\alpha=r$, whose left-hand side vanishes by
(H2), reduces to a single nonzero term).
\end{enumerate}
All sums converge absolutely (Cauchy--Schwarz on $\ell^2$ coefficients).
\end{proposition}

\begin{proof}
$g(\alpha\,\cdot)$ is almost periodic with Bohr--Fourier series
$\sum_\lambda a_\lambda e^{i(\alpha\lambda)x}$, since the Bohr mean is
invariant under positive dilations. Identity \eqref{eq:multAC} is the
two-function Parseval identity for Bohr (hence Besicovitch-$B^2$) almost
periodic functions~\cite{Besicovitch32,Corduneanu68} applied to the pair
$g(\alpha\,\cdot)$, $g(\beta\,\cdot)$. For \eqref{eq:addAC}: $|g|^2\equiv 1$
gives $M[|g|^2 e^{-i\tau x}] = 0$ for $\tau\neq0$, while the two-function
Parseval identity applied to $g$ and $e^{i\tau x}g$ evaluates the same
mean as 
\[
\sum_{\lambda-\mu=\tau} a_\lambda\overline{a_\mu}.
\]
 Absolute
convergence follows from Cauchy--Schwarz, the sums being inner products
of shifted, resp.\ rescaled-index, $\ell^2$ families. The difference
statement is immediate from \eqref{eq:addAC}. For (a) and (b),
hypothesis (H2) states that the left-hand side of \eqref{eq:multAC}
vanishes for all $\alpha\neq\beta$; the conclusions then follow as
indicated in the statement. That (H2) cannot be dropped is shown by
$g\equiv 1$: Bohr almost periodic and pointwise unimodular, yet
$\mathcal{F}=\{0\}$.
\end{proof}

\begin{remark}[Canonicity of the kernel $\sin(\sigma\gamma\log p)$]
Lemma~\ref{lem:canonicity} shows that, within the class of unimodular
trigonometric almost periodic kernels with pairwise orthogonal dilates,
the building block is forced to be the character $x\mapsto e^{icx}$
(up to constant phase and the sign of $c$); that much is rigid --- the
character part admits no alternative within the class. What follows is a
choice, and it is recorded as one: of the
real functions spanned by the $\pm c$-pair, the even one does not vanish at the
origin and the odd one does, so the convention that a mode contributes nothing
at $\log p=0$ selects $\sin(cx)$ up to a scalar. The substitutions
$c=\sigma\gamma_k$, $x=\log p$ then produce exactly the kernel
$\sin(\sigma\gamma_k\log p)$ used throughout the series. Class definitions
that vary the input data therefore run against a kernel whose character part
is rigid, with projection and frequency conventions declared once and held
fixed.
\end{remark}

For infinite spectra the conclusion of Lemma~\ref{lem:canonicity} is open.
The analytic ingredients are proved above --- \eqref{eq:multAC} for all Bohr
functions, \eqref{eq:addAC} under (H1), and the conclusions of
Proposition~\ref{prop:resonance} under (H1)--(H2) --- so what is missing is
purely combinatorial; the precise question is recorded in
\S\,\ref{sec:open}. The trigonometric class suffices here, since the
phasor kernel has finite spectrum. Beyond the single citation of the
two-function Parseval identity in the proof of
Proposition~\ref{prop:resonance}, no almost periodic machinery is used
anywhere in this series.

\subsection{Completion of the mode set}

\begin{remark}[Restriction to prime modes as a design choice; numerical]
\label{rem:lambda}
The coupling operator is built from unweighted prime phasors
$\sin(\sigma\gamma_k\log p)$. Completing the mode set to all prime powers
$n=p^m\le\kappa$, which gives $24$ modes at $\kappa=53$, leaves every
structural conclusion intact. The Gram identity of Theorem~\ref{thm:kernel}
holds verbatim for the extended label set: its proof uses only
$\sin a\sin b=\tfrac12[\cos(a-b)-\cos(a+b)]$ and the logarithmic labels, never
the primality of the label, so
$T_{m,n}=\tfrac12\{G(\log(m/n))-G(\log(mn))\}$ for the completed modes as
well --- an algebraic identity, with the residual $0$ at $30$ digits serving as
a reproduction check rather than as its ground. Further,
the curvature stays sign-definite negative across
$\varepsilon\in[0.04,0.07]$; and the spectral profile is unchanged, the
leading mode remaining the one associated with $\gamma_1$, while
\[
\lambda_{\max}/\mathrm{Tr}\quad\text{and}\quad\lambda_2/\lambda_1
\]
shift by under $4\%$.
The numerical effect of the completion on the curvature at
$\varepsilon=0.05$ is small and depends on the --- genuinely underdetermined
--- power-mode weight: $+3.8\%$ for the weight $\Lambda(n)\log n$, $+1.3\%$
for $\Lambda(n)^2$, and $+10.8\%$ for the operator's own second-moment weight
$(\log n)^2$. The completion does \emph{not} reproduce the Guinand--Weil prime
side: the raw phasors carry no functional-equation factor $n^{-1/2}$, whereas
the explicit-formula prime term weights each prime power by
$\Lambda(n)/\sqrt{n}$~\cite{IwaniecKowalski}. The restriction to primes is therefore a modelling convention with
quantified robustness: the curvature shifts are sign-stable and lie between
$1.3\%$ and $10.8\%$ depending on the power-mode weight; separately, the two
quoted spectral ratios move by less than $4\%$. It is not a truncation of a
canonical prime-power sum.
\end{remark}

% ============================================================
\section{Scaling Laws of the Renormalised Spectrum}
\label{sec:scaling}

Write
\[
r(x):=\frac{\lambda_{\max}(\Tt(x))}{\mathrm{Tr}\,\Tt(x)},\qquad u:=\log x,
\]
for the normalised leading spectral ratio at prime cutoff $x$.

\begin{theorem}[Non-convergence of the normalised leading ratio; proved,
separation certified]
\label{thm:nonconv}
At $\sigma=\tfrac12$ and for the reference ordinate and regularisation data
$(N,\varepsilon)=(100,0.05)$, the limit $\lim_{x\to\infty}r(x)$ does not
exist. Quantitatively,
\[
\limsup_{x\to\infty}r(x)-\liminf_{x\to\infty}r(x)\;\ge\;\Delta_0,
\qquad \Delta_0:=0.0840112 .
\]
\end{theorem}

\emph{Strategy of the proof.} Four steps: a uniform elementary asymptotic for
prime sums of $p^{i\theta}$ over the finite frequency set involved; entrywise
convergence of $\Tt(x)/\pi(x)$ to an explicit quasi-periodic limit matrix
$Q(u)$; transfer to eigenvalues by Weyl's inequality, which needs no spectral
gap; and a recurrence argument producing two values of $u$ that are approached
infinitely often, at which the limiting ratio
$R(u)=\lambda_{\max}(Q(u))/\mathrm{Tr}\,Q(u)$ is certified to differ.

\begin{proof}
\textit{Step 1.} With 
\[
h(\theta,u):=\Re[e^{i\theta u}/(1+i\theta)]
=(\cos\theta u+\theta\sin\theta u)/(1+\theta^2),
\]
 partial summation from the
prime number theorem with classical error term gives
\[
\frac{1}{\pi(x)}\sum_{p\le x}\cos(\theta\log p)=h(\theta,\log x)+O(1/\log x),
\]
uniformly over the finitely many frequencies occurring below; since
$|1+i\theta|\ge 1$ there is no small denominator, so no lower bound on
$|\theta|$ is needed.

\textit{Step 2.} Apply this entrywise with
$w_k=e^{-\varepsilon^2\gamma_k^2/2}$ to
\[
\Tt_{kl}(x)=w_kw_l\sum_{p\le x}\sin(\sigma\gamma_k\log p)\,
\sin(\sigma\gamma_l\log p),\qquad \sigma=\tfrac12 .
\]
The product formula used in
Theorem~\ref{thm:kernel} turns each entry into a difference of two cosine
sums, and yields the real symmetric limit matrix
\[
Q_{kk}(u)=\tfrac{w_k^2}{2}\bigl(1-h(\gamma_k,u)\bigr),
\qquad
Q_{kl}(u)=\tfrac{w_kw_l}{2}
\Bigl(h\bigl(\tfrac{\gamma_k-\gamma_l}{2},u\bigr)
     -h\bigl(\tfrac{\gamma_k+\gamma_l}{2},u\bigr)\Bigr)\ (k\neq l),
\]
with $\|\Tt(x)/\pi(x)-Q(\log x)\|\le K/\log x$ for an explicit $K$. The map
$u\mapsto Q(u)$ is quasi-periodic: it is a continuous function along a
Kronecker flow on a torus whose dimension is the number of rationally
independent frequencies among the $\gamma_k$ and their half-sums and
half-differences.

\textit{Step 3.} Since $|h(\theta,u)|\le(1+\theta^2)^{-1/2}$ for all $u$, the
trace is bounded below uniformly,

\[
\mathrm{Tr}\,Q(u)\ge\tfrac{\WeN}{2}\bigl(1-(1+\gamma_1^2)^{-1/2}\bigr)
=\tau_0=0.6380807\ldots>0
\]
 with
\[
\WeN:=\sum_{k=1}^{N}e^{-\varepsilon^2\gamma_k^2}=\sum_k w_k^2=1.37306
\]
at the reference data, and
$\|Q(u)\|\le\|Q(u)\|_F\le \WeN$, so $R(u)$ is well defined, bounded by
$\WeN/\tau_0\le 2.152$, and Lipschitz in $u$ with an explicit constant.
Weyl's inequality transfers Step~2 to $\lambda_{\max}$ without any gap
assumption, giving $|r(x)-R(\log x)|\le K'/\log x$.

\textit{Step 4.} By Dirichlet's pigeonhole principle the flow returns: there
are $\tau_n\to\infty$ with all relevant frequencies simultaneously close to
integer multiples of $2\pi$, hence $\|Q(u+\tau_n)-Q(u)\|\to 0$ uniformly in
$u$. Therefore every value of $R$ is a limit point of $r$. At the two explicit
points $u_1=17.6170$ and $u_2=7.1815$ the certificate below gives
$R(u_1)\ge 0.5169091132$ and $R(u_2)\le 0.4328979030$, whence the stated
separation. Finally $r$ is a step function, constant between consecutive
primes, so its value set over real cutoffs coincides with that over prime
cutoffs, and $\limsup$ and $\liminf$ agree in the two readings.
\end{proof}

\emph{Certification.} The separation of Step~4 is certified by ball
arithmetic, not by floating-point evaluation. The lower bound on
$\lambda_{\max}(Q(u_1))$ is the Rayleigh quotient $\langle v,Qv\rangle/
\langle v,v\rangle$ of an explicit vector, which is rigorous for any $v$
whatever; the upper bound on $\lambda_{\max}(Q(u_2))$ is
$(\mathrm{Tr}\,Q^{m})^{1/m}$ at $m=64$, which majorises the largest eigenvalue
for even $m$ without any positivity assumption; both traces are enclosed
directly and the two ratios are formed by outward-rounded interval division.
The ordinates enter as certified enclosures --- Arb's isolation of the first
hundred zeta zeros, whose ball midpoints are cross-checked against the
normative double-precision layer --- so the enclosures are themselves part of
the certificate rather than a hypothesis about it, and their radii are
propagated through every operation. The certificate returns
\[
R(u_1)-R(u_2)\;\ge\;0.0840112102,
\]
and is produced by \texttt{cert\_paper9\_ratio.py} (\S\,\ref{sec:methods}).

\begin{definition}[Averaging operator and ordinate profile]
\label{def:mean}
For locally integrable $\psi$ on $[2,\infty)$ set
\[
\langle\psi\rangle_X:=\frac{1}{\log X}\int_2^X\psi(x)\,\frac{dx}{x},
\qquad
\langle\psi\rangle:=\lim_{X\to\infty}\langle\psi\rangle_X
\]
whenever the limit exists. This logarithmic window mean is the only averaging
operator used in this paper; no other normalisation appears. The
\emph{ordinate profile} is
\[
f(x):=\sum_{k=1}^{N}e^{-\varepsilon^2\gamma_k^2}
\Bigl(\Im\frac{x^{i\gamma_k}}{\tfrac12+i\gamma_k}\Bigr)^{\!2}.
\]
\end{definition}

\begin{lemma}[Spectrum of the ordinate profile and its leading profile;
proved]
\label{lem:fmean}
Write $\tfrac12+i\gamma_k=\varrho_ke^{i\phi_k}$ with
$\varrho_k^2=\tfrac14+\gamma_k^2$, and set

\[
b_k:=e^{-\varepsilon^2\gamma_k^2}/(2\varrho_k^2).
\]
 Then $f$ is, in the
variable $u=\log x$, an exponential polynomial with frequency set
$\{0\}\cup\{\pm2\gamma_k\}$,
\[
f(e^u)=\langle f\rangle-\sum_{k=1}^{N}b_k\cos(2\gamma_ku-2\phi_k),
\qquad
\langle f\rangle=\sum_{k=1}^{N}b_k
=\tfrac12\sum_{k=1}^{N}\frac{e^{-\varepsilon^2\gamma_k^2}}
{\tfrac14+\gamma_k^2}>0 ,
\]
and $\langle f\rangle=0.0021613227$ at the reference parameters. The
associated \emph{leading profile}
\[
F(U):=\sum_{\nu}\frac{c_\nu e^{i\nu U}}{1+i\nu}
=\langle f\rangle-\sum_{k=1}^{N}b_k\,
\Re\frac{e^{i(2\gamma_kU-2\phi_k)}}{1+2i\gamma_k},
\]
where $f(e^u)=\sum_\nu c_\nu e^{i\nu u}$, is Bohr almost periodic with Bohr
mean $\langle F\rangle=c_0=\langle f\rangle$ and satisfies
\[
|F(U)-\langle f\rangle|\;\le\;
\sum_{k=1}^{N}\frac{b_k}{\sqrt{1+4\gamma_k^2}}
\;\le\;\frac{\langle f\rangle}{\sqrt{1+4\gamma_1^2}},
\]
\[
\text{hence}\qquad
F(U)\;\ge\;\langle f\rangle
\Bigl(1-\frac{1}{\sqrt{1+4\gamma_1^2}}\Bigr)\;>\;0
\]
for all $U$: the second inequality uses only $b_k>0$ and
$\gamma_k\ge\gamma_1>0$, so $F$ is bounded away from zero with no numerical
input. At the reference parameters the sharper half-width evaluates to
$\sum_{k}b_k/\sqrt{1+4\gamma_k^2}=6.7448532\cdot10^{-5}$, giving the
band $F(U)\in[\,0.0020938741,\;0.0022287713\,]$, and the exact envelope above
evaluates to $0.0020849161$; these three decimals are finite evaluations at the
reference parameters and carry no proof step [numerical].
\end{lemma}

\begin{proof}
\[
\Im\Bigl[\frac{x^{i\gamma_k}}{\tfrac12+i\gamma_k}\Bigr]
=\varrho_k^{-1}\sin(\gamma_ku-\phi_k),
\]
and $\sin^2\vartheta=\tfrac12(1-\cos2\vartheta)$ gives the displayed expansion
with $b_k$ as stated; the frequency set is $\{0\}\cup\{\pm2\gamma_k\}$ and the
sum is finite, so no convergence question arises. The Bohr mean of
$e^{i\nu U}$ vanishes for $\nu\neq0$, whence $\langle F\rangle=c_0$, and
$c_0=\langle f\rangle$ because dividing by $1+i\nu$ leaves the constant mode
unchanged. The first bound is the triangle inequality applied to the displayed sum,
using $|1+2i\gamma_k|=\sqrt{1+4\gamma_k^2}$. The second follows termwise from
$\gamma_k\ge\gamma_1>0$ together with $b_k>0$, and summing gives
$\sum_k b_k=\langle f\rangle$; since $\gamma_1>0$ the factor
$1-(1+4\gamma_1^2)^{-1/2}$ is positive, so the envelope is strictly positive for
every $N$ and every $\varepsilon>0$. The decimals are finite evaluations at the
reference parameters and enter no step of the argument.
\end{proof}

\begin{theorem}[Divergence of the energy-asymmetry functional in mean; proved]
\label{thm:eta}
Let $c_p^{\eta}=\sqrt{V_p(\tfrac12)}$ and let the phasor argument be
$\gamma_k\log p$, as in the paper of the series in which
$\eta_{\mathrm{orig}}$ was introduced. Then, with $F$ the leading profile of
Lemma~\ref{lem:fmean},
\[
\langle\eta_{\mathrm{orig}}\rangle_X
=-\frac{4}{\WeN}\,F(\log X)\,\frac{X}{(\log X)^3}
\;+\;O\!\left(\frac{X}{(\log X)^4}\right)
\;\longrightarrow\;-\infty .
\]
The leading term carries a bounded oscillating profile, not a constant: the
window mean is dominated by the upper endpoint of the window, so the
oscillation of $f$ survives the averaging instead of being washed out. What
makes the divergence unconditional is the exact envelope of
Lemma~\ref{lem:fmean},
$F\ge\langle f\rangle\bigl(1-(1+4\gamma_1^2)^{-1/2}\bigr)>0$, together with
$\langle F\rangle=\langle f\rangle$.
\end{theorem}

\begin{proof}
\textit{Step 1 (structural energy).} With $a_p$ the $p$-th column,

\[
E_{\mathrm{str}}(x)=\sum_{p\le x}(c_p^{\eta})^2\|a_p\|^2,
\]
 where
\[
(c_p^{\eta})^2=\frac{4(\log p)^2p}{(p-1)^2}
=\frac{4(\log p)^2}{p}\bigl(1+O(p^{-1})\bigr),
\qquad
\|a_p\|^2=\frac{\WeN}{2}-\frac12\sum_kw_k^2\cos(2\gamma_k\log p).
\]
Mertens'
theorem gives 
\[
\sum_{p\le x}(\log p)^2/p=\tfrac12(\log x)^2+O(\log x),
\]
 so the
constant part contributes $\WeN(\log x)^2(1+o(1))$. For the oscillating part, write $\vartheta(t)=t+E(t)$ with
$E(t)=O\bigl(te^{-c\sqrt{\log t}}\bigr)$, the prime number theorem with
classical error term~\cite{IwaniecKowalski}. For fixed
$\tau=2\gamma_k\neq0$,
\[
S_\tau(x):=\sum_{p\le x}(\log p)\,p^{-1+i\tau}
=\int_2^x t^{-1+i\tau}\,d\vartheta(t),
\]
\[
\int_2^x t^{-1+i\tau}\,d\vartheta(t)
=\frac{x^{i\tau}-2^{i\tau}}{i\tau}
+\Bigl[\frac{E(t)}{t^{1-i\tau}}\Bigr]_2^x
+(1-i\tau)\int_2^x\frac{E(t)}{t^{2-i\tau}}\,dt
=O_\tau(1),
\]
the boundary term vanishing as $t\to\infty$ and the error integral being
absolutely convergent because
$\int_2^\infty e^{-c\sqrt{\log t}}\,dt/t<\infty$ (substitute $u=\log t$).
A second partial summation,

\[
\sum_{p\le x}(\log p)^2p^{-1+i\tau}
=(\log x)S_\tau(x)-\int_2^xS_\tau(t)\,dt/t=O_\tau(\log x),
\]
 then bounds
\[
\sum_{p\le x}\frac{(\log p)^2}{p}\cos(2\gamma_k\log p)=O(\log x).
\]
Summing the finitely many $k$,
\[
E_{\mathrm{str}}(x)=\WeN(\log x)^2+O(\log x),
\]
exactly the precision Step~3 needs for its
$O\bigl(x/(\log x)^{3}\bigr)$ quotient error.

\textit{Step 2 (spectral energy).} Since
$c_p^{\eta}=2(\log p)p^{-1/2}(1+O(p^{-1}))$, partial summation from
$\vartheta(x)=x+O(x/(\log x)^{A})$ gives, for each fixed $\gamma$,
\[
\sum_{p\le x}(\log p)\,p^{-1/2+i\gamma}
=\frac{x^{1/2+i\gamma}}{\tfrac12+i\gamma}+O\!\bigl(x^{1/2}(\log x)^{-A+1}\bigr).
\]
Taking imaginary parts, the $k$-th component of
$\sum_{p\le x}c_p^{\eta}a_p$ equals

\[
2w_k\sqrt{x}\,\Im[x^{i\gamma_k}/(\tfrac12+i\gamma_k)]
\]
 up to the same error,
so
\[
E_{\mathrm{spec}}(x)=4x\,f(x)+O\bigl(x(\log x)^{-A+1}\bigr).
\]

\textit{Step 3 (quotient).} Hence
\[
\eta_{\mathrm{orig}}(x)=1-\frac{E_{\mathrm{spec}}(x)}{E_{\mathrm{str}}(x)}
=-\,\frac{4x\,f(x)}{\WeN(\log x)^{2}}
+O\!\left(\frac{x}{(\log x)^{3}}\right).
\]

\textit{Step 4 (the average is endpoint-dominated).} Substituting $x=e^u$ in
Definition~\ref{def:mean} turns the window mean into
\[
\langle\eta_{\mathrm{orig}}\rangle_X
=-\frac{4}{\WeN\log X}\int_{\log 2}^{\log X}f(e^u)\,\frac{e^u}{u^2}\,du
+O\!\left(\frac{X}{(\log X)^4}\right),
\]
an integral whose integrand grows like $e^u$ and which is therefore governed
by its upper endpoint, not by the mean of $f$. Expanding $f(e^u)$ by
Lemma~\ref{lem:fmean} and integrating each mode,
\[
\int^{U}\frac{e^{(1+i\nu)u}}{u^{2}}\,du
=\frac{e^{(1+i\nu)U}}{(1+i\nu)U^{2}}
+O\!\left(\frac{e^{U}}{U^{3}}\right),
\]
uniformly over the finitely many frequencies $\nu$, because
$|1+i\nu|\ge1$. Summing over the modes reproduces exactly the combination
defining $F$, so the integral equals $F(U)\,e^{U}U^{-2}+O(e^UU^{-3})$ with
$U=\log X$. Dividing by $\log X$ gives the statement. Since
$F\ge\langle f\rangle\bigl(1-(1+4\gamma_1^2)^{-1/2}\bigr)>0$ by
Lemma~\ref{lem:fmean}, the mean diverges to
$-\infty$; the constant-rate form obtained by replacing $F$ with its Bohr
mean $\langle f\rangle$ is \emph{not} valid pointwise in $X$, only on
average.
\end{proof}

\emph{Numerical observation} (sign change of the quotient). The theorem is a
statement about means and says nothing about individual cutoffs. Evaluating
$\eta_{\mathrm{orig}}(\kappa)$ directly at every prime cutoff
$\kappa\le 2\cdot10^{5}$, with the same ordinate list and the weights above,
the quotient is positive at $\kappa=146521$, where it equals $+0.0049$,
negative at the next prime cutoff $146527$, where it equals $-0.0005$, and
negative at every larger prime cutoff in the search range. For orientation,
$\eta_{\mathrm{orig}}(53)=0.69078176$ and
$\eta_{\mathrm{orig}}(1009)=0.810955$, while
$\eta_{\mathrm{orig}}(10^{6})=-41.307$. These are evaluations on a finite
grid, not part of the theorem.

\begin{theorem}[Ordinate tail and the true limit parameters; proved]
\label{thm:sat}
Let $T_N=\Phi_N^{*}\Phi_N$ on $\Hstr$ be the loop operator built from the
first $N$ ordinates at fixed $\sigma=\tfrac12$, and $T_\infty$ the
corresponding sum over all positive ordinates; assume
$\mathrm{Tr}\,T_N>0$. Then
\[
\|T_\infty-T_N\|\;\le\;\mathrm{Tr}(T_\infty-T_N)\;\le\;
\pi(\kappa)\sum_{k>N}e^{-\varepsilon^2\gamma_k^2},
\qquad
|r_\infty-r_N|\;\le\;\frac{2\pi(\kappa)}{\mathrm{Tr}\,T_N}
\sum_{k>N}e^{-\varepsilon^2\gamma_k^2},
\]
and the tail is bounded by a geometric series against the
Riemann--von-Mangoldt density, so $r_N\to r_\infty$: for fixed $\kappa$ and $\varepsilon>0$ the rate is
governed by the Gaussian ordinate tail $e^{-\varepsilon^2\gamma_N^2}$, up to
the Riemann--von-Mangoldt factor. The limit parameters of the construction are therefore
$(\kappa,\varepsilon)$; $N$ is a truncation, not a parameter of the limit.
\end{theorem}

\begin{proof}
\textit{Step 1 (positive semidefiniteness and trace tail).}
On $\Hstr$ the difference decomposes into rank-one terms,

\[
T_\infty-T_N=\sum_{k>N}e^{-\varepsilon^2\gamma_k^2}v_kv_k^{\!\top}
\]
 with
$(v_k)_p=\sin(\sigma\gamma_k\log p)$, all positive semidefinite; hence the
operator norm is at most the trace, and $\|v_k\|^2\le\pi(\kappa)$ gives the
first bound.

\textit{Step 2 (ratio perturbation).} For the second, Weyl's inequality gives

\[
|\lambda_{\max,\infty}-\lambda_{\max,N}|\le\|\Delta\|
\]
 with
$\Delta=T_\infty-T_N$, and
\[
|r_\infty-r_N|
=\Bigl|\frac{\lambda_{\infty}}{\mathrm{Tr}\,T_\infty}
      -\frac{\lambda_{N}}{\mathrm{Tr}\,T_N}\Bigr|
\le\frac{\|\Delta\|\,\mathrm{Tr}\,T_N+\lambda_N\,\mathrm{Tr}\,\Delta}
        {\mathrm{Tr}\,T_\infty\,\mathrm{Tr}\,T_N}
\le\frac{2\,\mathrm{Tr}\,\Delta}{\mathrm{Tr}\,T_N},
\]
using 
\[
\lambda_N\le\mathrm{Tr}\,T_N\le\mathrm{Tr}\,T_\infty.
\]


\textit{Step 3 (Riemann--von-Mangoldt Gaussian tail).} For the tail, split $\sum_{k>N}e^{-\varepsilon^2\gamma_k^2}$ into dyadic blocks
$\gamma\in[T,2T)$ with $T=\gamma_N$: the Riemann--von-Mangoldt density gives
$O(T\log T)$ ordinates per block, each contributing at most
$e^{-\varepsilon^2T^2}$, and the resulting series in $T$ is dominated by a
geometric one. A single-term estimate would not suffice and is not used.

\textit{Step 4 (transfer to the companion Gram operator).}
The statement transfers to $\Tt_N=\Phi_N\Phi_N^{*}$ without further work:
the non-zero spectra of $T_N$ and $\Tt_N$ coincide (\S\,\ref{sec:operator}),
so trace, largest eigenvalue and hence $r_N$ agree in the two readings.
\end{proof}

\emph{Numerical observation} (double-precision threshold). The tail bound is
$2.3\cdot10^{-23}$ at $N=50$ and $1.8\cdot10^{-61}$ at $N=100$, against an
actual $|r_{100}-r_{50}|=2.9\cdot10^{-24}$. The increments therefore fall
below the double-precision unit roundoff of $r$ well before the reference
truncation: the first $N$ with $|r_N-r_{N-1}|<5.5\cdot10^{-17}$ is $N=39$.
Beyond that point the dependence on $N$ is invisible in double precision but
is not zero: at $40$ digits the increments are still resolvable and strictly
signed, with
\[
r_{51}-r_{50}=-2.10567\cdot10^{-24},
\qquad
\mathrm{Tr}(T_{51}-T_{50})=+6.62880\cdot10^{-23},
\]
so exact independence of $N$ is false and only the limit statement above is
correct. These control values were computed independently at $40$ digits from
the normative ordinate list; the digits beyond the tenth are not quoted,
because the normative list carries double precision and does not support
them.

Two further facts locate the mechanism. \emph{Numerical observation}
(leading-mode localisation). The leading eigenvector is not carried by
diagonal dominance: the largest diagonal entry stays bounded as the cutoff
grows, while $\lambda_{\max}$ grows like $\pi(\kappa)$. The leading mode is
the wave associated with $\gamma_1$: at cutoff $\kappa\approx10^{4}$, with
the $650$-ordinate extended list of \S\,\ref{sec:methods} and
$\varepsilon=0.05$, the squared overlap $|\langle v_{\max},e_1\rangle|^{2}$
against the corresponding ordinate direction is $0.94$, and it remains above
$0.83$ on the tested cutoff ladder up to $\kappa\approx10^{7}$. Whether the oscillation of $R$ has a
pointwise, as opposed to a recurrent, limiting behaviour is a diophantine
question, recorded in \S\,\ref{sec:open}.

\section{Class Invariance over Beurling Systems}
\label{sec:class}

Nothing in \S\,\ref{sec:scaling} used the multiplicative structure of the
integers. This section replaces the primes by a Beurling generalised prime
system and shows that the same proofs run --- carrying over the limiting
function $R$, the certified separation $\Delta_0$, the leading profile $F$ and
its band unchanged.

\begin{definition}[The Beurling side]
\label{def:beurling}
A \emph{Beurling system} is a non-decreasing unbounded sequence
$\mathcal{P}_{\mathfrak{B}}=\{q_1\le q_2\le\cdots\}$ with $q_1>1$, with counting function
$\pi_{\mathfrak{B}}(x)=\#\{j:q_j\le x\}$ and Chebyshev function
$\vartheta_{\mathfrak{B}}(x)=\sum_{q_j\le x}\log q_j$~\cite{Beurling37,DiamondZhang16}. The
coupling data are transported verbatim: the map
\[
\Phi_{\mathfrak{B}}(\sigma)_{k,q}=e^{-\varepsilon^2\gamma_k^2/2}\sin(\sigma\gamma_k\log q),
\qquad q\in\mathcal{P}_{\mathfrak{B}},\ q\le x,
\]
the operator

\[
\Tt_{\mathfrak{B}}(x)=\Phi_{\mathfrak{B}}\Phi_{\mathfrak{B}}^{*},
\]
 the ratio

\[
r_{\mathfrak{B}}(x)=\lambda_{\max}(\Tt_{\mathfrak{B}}(x))/\mathrm{Tr}\,\Tt_{\mathfrak{B}}(x),
\]
 the weights
$c_q^{\eta}=2(\log q)\sqrt{q}/(q-1)$, the column vectors
\[
a_q^{\eta}:=\bigl(e^{-\varepsilon^2\gamma_k^2/2}\sin(\gamma_k\log q)\bigr)_{k=1}^{N},
\]
the structural and spectral energies
\[
E^{\mathfrak{B}}_{\mathrm{str}}(x)=\sum_{q\le x}(c_q^{\eta})^2\|a_q^{\eta}\|^2,
\qquad
E^{\mathfrak{B}}_{\mathrm{spec}}(x)=\Bigl\|\sum_{q\le x}c_q^{\eta}\,a_q^{\eta}\Bigr\|^2,
\]
 and the quotient

\[
\eta_{\mathfrak{B}}=1-E^{\mathfrak{B}}_{\mathrm{spec}}/E^{\mathfrak{B}}_{\mathrm{str}}.
\]
 Two conventions coexist
here exactly as in the classical case: $r_{\mathfrak{B}}$ uses $\Phi_{\mathfrak{B}}(\tfrac12)$, while
$\eta_{\mathfrak{B}}$ keeps the no-$\sigma$ phasor convention of \S\,\ref{sec:scaling},
matching the import of $\eta_{\mathrm{orig}}$. The ordinate side is not
varied here: the $\gamma_k$ remain the ordinates of the zeta function.
For $\delta>0$, $C>0$ and $x_0\ge 2$ let
\[
\Dclass(C,x_0):=\bigl\{\mathfrak{B}:\ |\vartheta_{\mathfrak{B}}(x)-x|\le C\,x(\log x)^{-2-\delta}
\ \text{ for all } x\ge x_0\bigr\},
\qquad
\Dclass:=\bigcup_{C,x_0}\Dclass(C,x_0).
\]
For the uniformity clause write
\[
\Dclass(C,x_0,M_0):=\bigl\{\mathfrak{B}\in\Dclass(C,x_0):\ \pi_{\mathfrak{B}}(x_0)\le M_0\bigr\}:
\]
a bound on the early data is genuinely additional information, because points
accumulating just above $1$ enter $\pi_{\mathfrak{B}}$ at full strength while contributing
almost nothing to $\vartheta_{\mathfrak{B}}$.
\end{definition}

\begin{theorem}[Class transfer: same limiting objects and same profile;
proved]
\label{thm:class}
Keep the same reference ordinate and regularisation data as in
Theorem~\ref{thm:nonconv}, and vary only the prime-side Beurling system.
Let $\mathfrak{B}\in\Dclass$. Then 
\[
r_{\mathfrak{B}}(x)=R(\log x)+O_{\mathfrak{B}}(1/\log x)
\]
 with the \emph{same}
limiting function $R$ as in Theorem~\ref{thm:nonconv}; the limit fails to
exist for every such $\mathfrak{B}$; and

\[
\limsup r_{\mathfrak{B}}-\liminf r_{\mathfrak{B}}\ge\Delta_0
\]
 with the same $\Delta_0=0.0840112$.
Likewise $\langle\eta_{\mathfrak{B}}\rangle_X\to-\infty$ for every $\mathfrak{B}\in\Dclass$ with the
same leading profile $F$ and the same exact positivity envelope of
Lemma~\ref{lem:fmean} in the rate of Theorem~\ref{thm:eta},
class-independently. The objects $R$, $u_1$, $u_2$, $\Delta_0$, $F$ and its
Bohr mean $\langle f\rangle$ depend only on the ordinate side and are free of
$\mathfrak{B}$; the certificate of Theorem~\ref{thm:nonconv} therefore transfers
verbatim.
Uniformity holds on each $\Dclass(C,x_0,M_0)$ with the implied constants
depending only on $(C,x_0,\delta,M_0)$; over the full class $\Dclass$ it
fails.
\end{theorem}

\begin{proof}
\textit{Step 1.} The only property of the primes used in Step~1 of
Theorem~\ref{thm:nonconv} is $\vartheta(x)=x+o(x/\log x)$, which
$\Dclass$ supplies by hypothesis. Split the sum at $x_0$: the early points
contribute at most $M_0$ to the sum directly and at most $M_0$ again through
the boundary term of the partial summation at $x_0$ --- together at most
$2M_0/\pi_{\mathfrak{B}}(x)$ after normalisation --- and
$\pi_{\mathfrak{B}}(x)\ge(1+o(1))\,x/\log x$ follows from the $\vartheta$-condition
itself, so the early-data term is absorbed in the error. For $q>x_0$ partial
summation against the $\vartheta$-bound reproduces
\[
\frac{1}{\pi_{\mathfrak{B}}(x)}\sum_{q\le x}\cos(\theta\log q)=h(\theta,\log x)+O(1/\log x)
\]
with the same $h$ and the implied constant depending only on
$(C,x_0,\delta,M_0)$. Steps~2--4 use only this asymptotic and the fixed ordinate
data, so they carry over with the same $Q$, the same $R$, and the same
certified points.

\textit{Step 2.} For the quotient, the Beurling form of Mertens' theorem

\[
\sum_{q\le x}(\log q)^2/q=\tfrac12(\log x)^2+O_{\mathfrak{B}}(1)
\]
 follows from
Definition~\ref{def:beurling} by partial summation, the error integral
$\int^\infty(\log t)^{-2-\delta}\,dt/t$ converging precisely because
$\delta>0$; and the spectral side is again $4xf(x)$ up to a smaller error,
$f$ being built from the ordinates alone. The constants $\WeN$ and the whole profile $F$, band included, are therefore
untouched.

\textit{Step 3 (uniformity).} Within $\Dclass(C,x_0,M_0)$ every implied
constant above is a function of $(C,x_0,\delta,M_0)$ only. Over the full
class $\Dclass$, uniformity fails, and this is a separate statement about
the union, established by finite deletion: deleting finitely many elements
from $\mathcal{P}_{\mathfrak{B}}$ keeps $\mathfrak{B}$ in $\Dclass$ but moves $x_0$ arbitrarily far out, so no
threshold can be chosen uniformly.
\end{proof}

\begin{remark}[Finite insertion near $1$]
Fix $(C,x_0,\delta)$ and a base system in $\Dclass(C,x_0,M_0)$ with fixed
positive uniform slack: there exists $C'<C$ with

\[
|\vartheta_{\mathfrak{B}}(x)-x|\le C'\,x(\log x)^{-2-\delta}
\]
 for all
$x\ge x_0$. Set $m_0:=\inf_{x\ge x_0}x(\log x)^{-2-\delta}>0$ and choose
$0<\varepsilon_0\le(C-C')\,m_0$. Adjoining $K$ generalized primes at
$q_K=\exp(\varepsilon_0/K)<x_0$ costs only $\varepsilon_0$ in $\vartheta$,
independently of $K$, while raising the count at $x_0$ by $K$; with
$b(x):=x(\log x)^{-2-\delta}$ the perturbed system stays in the class,
since 
\[
|E_{\mathrm{new}}(x)|\le C'\,b(x)+\varepsilon_0\le C\,b(x).
\]
 The
normalized cosine averages used in Step~1 of the proof are then dominated
by the inserted points, so no equidistribution-based error control
depending only on $(C,x_0,\delta)$ can be uniform over $\Dclass(C,x_0)$.
The observables themselves are far less sensitive: the inserted points
shift $r(x)$ by $O(1/K)$ and the windowed $\eta$-mean by an amount bounded
in terms of the reference data alone, uniformly in $K$, within the error
budget of Theorem~\ref{thm:class}. The construction therefore obstructs
any proof route via uniform equidistribution control; it does not decide
whether the conclusions of Theorem~\ref{thm:class} themselves persist on
$\Dclass(C,x_0)$.
\end{remark}

The exponent marks the majorant threshold of the present argument, and
nothing stronger is claimed. Writing $E=\vartheta_{\mathfrak{B}}(t)-t$ and summing by
parts,
\[
\sum_{q\le x}\frac{(\log q)^2}{q}
=\tfrac12(\log x)^2+\Bigl[\frac{E(t)\log t}{t}\Bigr]_2^x
+\int_2^x E(t)\,\frac{\log t-1}{t^2}\,dt ,
\]
the boundary term is $O((\log x)^{-1-\delta})$ and the integral is majorised
by $\int^x dt/(t(\log t)^{1+\delta})$. For $\delta>0$ this majorant converges
and the structural energy retains a bounded remainder; at $\delta=0$ it grows
like $\log\log x$; for $\delta<0$ it is $O((\log x)^{-\delta})$. These are
upper bounds delivered by this argument, not matching lower bounds: whether
some system at $\delta\le0$ nevertheless keeps a bounded remainder is part of
the boundary question of Open Problem~\ref{op:order}. The direction of the classical implication should be stated
carefully: a sharp Mertens estimate implies the prime number theorem always,
whereas the prime number theorem does not imply a sharp Mertens
estimate~\cite{DebruyneVindas17}; counterexamples exist.

The next statement concerns a different object, and the difference matters.
It is not a Beurling system: it is a \emph{smooth generalised-prime measure},
a continuous density with a bounded oscillation, and its role here is to show
how the constants respond to a perturbation of the density that is not a
change of class.

\begin{definition}[Smooth generalised-prime measure and the conjugate profile]
\label{def:smoothosc}
For $0\le a<1$ and $\omega>0$ let

\[
d\vartheta_{a,\omega}(t)=\bigl(1+a\sin(\omega\log t)\bigr)\,dt,
\]
 and write
\[
\tilde h(\varphi,u):=\Im\frac{e^{i\varphi u}}{1+i\varphi}
=\frac{\sin\varphi u-\varphi\cos\varphi u}{1+\varphi^{2}}
\]
for the conjugate of the density profile $h$ of
Theorem~\ref{thm:nonconv}. The sine convention is the one in which the
transfer below takes its stated form; with a cosine density the same
computation gives $h(\theta+\omega,u)+h(\theta-\omega,u)$ in place of the
difference of conjugates, and the mass factor changes accordingly. The two
conventions are not interchangeable term by term.
\end{definition}

\begin{theorem}[Sideband transfer of the leading representation; proved
(smooth family), numerical in the discrete embedding]
\label{thm:sideband}
Assume the non-resonance condition: the frequencies $\gamma_k$ and
$\gamma_k\pm\omega$, $k\le N$, are all non-zero and pairwise distinct in
absolute value. Then for $d\vartheta_{a,\omega}$ the leading representation transfers
with explicit replacements:
\[
h(\theta,u)\;\longmapsto\;\hat h(\theta,u)
:=h(\theta,u)+\frac{a}{2}\bigl[\tilde h(\theta+\omega,u)
-\tilde h(\theta-\omega,u)\bigr],
\]
the limit matrix being formed from $\hat h$ by the element formulas of
Theorem~\ref{thm:nonconv} with $\hat h(0,u)=1+a\,\tilde h(\omega,u)$ on the
diagonal, and the Bohr mean of the leading profile shifting to
\[
\langle f_{\mathrm{osc}}\rangle=\langle f\rangle+\frac{a^2}{8}
\sum_{k=1}^{N}e^{-\varepsilon^2\gamma_k^2}
\left[\frac{1}{\tfrac14+(\gamma_k+\omega)^2}
     +\frac{1}{\tfrac14+(\gamma_k-\omega)^2}\right].
\]
A bounded relative oscillation cannot change the \emph{order} of the spectral
energy, by linearity in the measure; it can only move these objects.
No general assertion is made that the limiting ratio formed from $\hat h$ is
non-constant, or that $\inf_U F_{\mathrm{osc}}(U)>0$, for all admissible
$(a,\omega)$; for the reference pair $(a,\omega)=(0.3,2\pi/3)$ the
numerical embedding stands as stated ($0.25\%$).
\end{theorem}

\begin{proof}
\textit{Count-side transfer.} Insert $d\vartheta_{a,\omega}$ into the
normalised counting argument of Theorem~\ref{thm:nonconv}, Step~1. Writing

\[
\sin(\omega\log t)=(t^{i\omega}-t^{-i\omega})/2i,
\]
introduce the count measure
\[
d\Pi_{a,\omega}(t):=\frac{d\vartheta_{a,\omega}(t)}{\log t}
=\frac{1+a\sin(\omega\log t)}{\log t}\,dt,
\]
so that
\[
\int_2^x t^{i\theta}\,d\Pi_{a,\omega}(t)
=\frac{x^{1+i\theta}}{(1+i\theta)\log x}
+\frac{a}{2i}\Bigl[
\frac{x^{1+i(\theta+\omega)}}{(1+i(\theta+\omega))\log x}
-\frac{x^{1+i(\theta-\omega)}}{(1+i(\theta-\omega))\log x}
\Bigr]
+O\!\Bigl(\frac{x}{\log^2x}\Bigr).
\]
After multiplication by $\log x/x$ and taking real parts this is exactly
the stated $\widehat h$ in the reference normalization; if the perturbed
mass itself is used for normalization, divide by the corresponding
$\theta=0$ term.

\textit{Spectral-energy transfer.} Insert $d\vartheta_{a,\omega}$ into the
partial summation of Theorem~\ref{thm:eta}, Step~2:
\[
\int_2^x t^{-1/2+i\gamma}\,d\vartheta_{a,\omega}(t)
=\frac{x^{1/2+i\gamma}}{\tfrac12+i\gamma}
+\frac{a}{2i}\Bigl[
\frac{x^{1/2+i(\gamma+\omega)}}{\tfrac12+i(\gamma+\omega)}
-\frac{x^{1/2+i(\gamma-\omega)}}{\tfrac12+i(\gamma-\omega)}
\Bigr]+O(1).
\]
Under the theorem's non-resonance hypothesis the Bohr cross terms vanish,
giving the printed $a^{2}/8$ shift: the two sidebands contribute
\[
\frac{a^2}{8}\Bigl[\bigl|\tfrac12+i(\gamma_k+\omega)\bigr|^{-2}
+\bigl|\tfrac12+i(\gamma_k-\omega)\bigr|^{-2}\Bigr].
\]
Since the perturbation is a bounded
multiplicative factor of the measure, the leading order $x$ of the spectral
energy and the order $(\log x)^2$ of the structural energy are unchanged.
\end{proof}

\begin{remark}[Exact resonance]
At $\omega=\gamma_j$ the sideband of the $j$-th mode lands at frequency zero
and becomes a constant, which the Bohr mean no longer halves: the mean
acquires the additional term
\[
\tfrac{a^2}{2}\sum_{k:\,\gamma_k=\omega}e^{-\varepsilon^2\gamma_k^2}
\]
beyond the displayed formula --- twice the generic contribution of that
sideband. The registered probe $(a,\omega)=(0.3,2\pi/3)$ satisfies the
non-resonance condition, so all figures quoted below are unaffected.
\end{remark}

\begin{remark}[Which quantity is being shifted]
The displayed shift of the mean concerns the unnormalised reading. If one
normalises by the perturbed total mass, the same computation carries a factor
$\bigl(1+a\,\tilde h(\omega,u)\bigr)^{-1}$, and the shift is correspondingly
smaller; the two readings differ at order $a^2$ and must not be mixed. The
$a^2$ statement above is a statement about the Bohr-mean part of the profile;
the window behaviour runs through the perturbed profile $F_{\mathrm{osc}}$ and
is not obtained by substituting a constant. All figures quoted here are
unnormalised.
\end{remark}

\emph{Numerical evidence} (embedding). For the probe
$(a,\omega)=(0.3,2\pi/3)$ the formula predicts
$\langle f_{\mathrm{osc}}\rangle/\langle f\rangle=1.04749$. An explicitly
constructed discrete test world realising that density profile gives
$1.0449$, a deviation of $0.25\%$, and the predicted limit matrix matches the
measured leading ratio of that world to $7\cdot10^{-4}$. The closed formula is
proved for the smooth family; the agreement of a discrete realisation with it
is numerical, and is labelled as such.

What the two scaling laws provably see of the prime side, at leading
order, is the density class; the remainder terms may retain
system-dependent information. At that leading order the limiting function
$R$, the separation $\Delta_0$ and the profile $F$ take the same values
across $\Dclass$. The statement is about
these observables and is not a claim that no observable whatever could
distinguish two systems within the class. Two boundary facts complete the picture. Degrading the error exponent to
$x(\log x)^{-\alpha}$ preserves the laws for \emph{every} $\alpha>0$, with
degraded rates. Indeed, writing 
\[
E(t)=\vartheta_{\mathfrak{B}}(t)-t=O(t/(\log t)^{\alpha}),
\]

direct Stieltjes partial summation against $E$ gives
\[
\int_{2}^{x}t^{-1/2+i\gamma}\,dE(t)
=O_{\alpha,\gamma}\!\Bigl(\frac{\sqrt{x}}{(\log x)^{\alpha}}\Bigr),
\]
the boundary term and the integral against $|E(t)|\,t^{-3/2}$ both being of
this order (split the integral at $\sqrt{x}$), so the spectral energy keeps
its leading term:

\[
E^{\mathfrak{B}}_{\mathrm{spec}}(x)=4x\,f(x)+O_{\alpha,N,\varepsilon}\bigl(x/(\log
x)^{\alpha}\bigr)
\]
 --- small against $x$ for every $\alpha>0$, where the
substitution $A=\alpha$ in the proof of Theorem~\ref{thm:eta} would require
$\alpha>1$. On the structural side the remainder follows the
$\delta$-case distinction recorded above under $\alpha=2+\delta$: bounded
for $\alpha>2$, $\log\log x$ at $\alpha=2$, $O((\log x)^{2-\alpha})$
below --- rates degrade, the leading order never does. And it is not settled here whether the window bound admits an exceptional-set
control under $\Dclass$ at all: the tail available to the present argument is
polynomial rather than summable, which would cap any such control at density
tending to zero, and the question is recorded in \S\,\ref{sec:open}. Whether a genuine order deviation destroys the laws is open
(\S\,\ref{sec:open}).

\section{The Certified Factorisation Barrier}
\label{sec:barrier}

Section~\ref{sec:class} shows that, at the level of the proved leading
laws, variation within $\Dclass$ leaves $R$, $\Delta_0$, $F$, and the
stated band unchanged; lower-order discrimination and discrimination by
other observables are not excluded. This section asks what the stated
curvature observable can distinguish against control objects outside that
within-class comparison, and certifies the margin. Three control objects appear and are
kept strictly apart.

\begin{definition}[The three control objects]
\label{def:controls}
Let $\mu_{\mathrm{sm}}^{\infty}$ be the smooth measure $dt/\log t$ on
$[2,\infty)$, and let $\mu_{\mathrm{sm}}$ be its restriction to $[2,\kappa]$,
of total mass $M=\int_2^{\kappa}dt/\log t=18.1846895$ at the reference
cutoff; every finite comparison below runs exclusively over the
restriction.
Let $\mu_{\mathrm{norm}}:=(16/M)\,\mu_{\mathrm{sm}}$ be its rescaling to the
mass of the prime counting measure, so that $\mu_{\mathrm{norm}}$ --- and only
$\mu_{\mathrm{norm}}$ --- matches the prime count exactly. Let $A$ be the
discrete quadrature world --- from here on, $A$ denotes exclusively the discrete
quadrature world; the Beurling system index in \S\,\ref{sec:class} is
$\mathfrak{B}$ --- given by the $16$ nodes $a_1<\cdots<a_{16}$ in $[2,\kappa]$
determined by
\[
\int_2^{a_j}\frac{dt}{\log t}=\Bigl(j-\tfrac12\Bigr)\frac{M}{16},
\qquad j=1,\dots,16,
\]
that is the midpoints of the $16$ cells of equal $\mu_{\mathrm{sm}}$-mass
$M/16$, carrying unit weights. Write
$E_\pi(t):=\#\{p\le t\}-\mu_{\mathrm{sm}}([2,t])$ for the resulting
\emph{counting} error of $\pi$ against the smooth measure. Curvature
sums attached to these objects are written $\Bsm$, $B[\mu_{\mathrm{norm}}]$
and $B[A]$; the prime curvature is $B[P]$.
\end{definition}

\begin{remark}[The smooth measure carries no zeros]
\label{rem:smoothzeta}
The zeta function attached to the \emph{infinite} measure
$\mu_{\mathrm{sm}}^{\infty}$ is

\[
\zeta_{\mathrm{sm}}(s)=\exp\bigl(E_1((s-1)\log 2)\bigr)
\]
 for $\Re s>1$, continued
analytically to the slit plane; an exponential never vanishes, so
$\zeta_{\mathrm{sm}}$ is zero-free, and its singular behaviour at $s=1$ is the one the density
$dt/\log t$ requires. The transform of the finite restriction
$\mu_{\mathrm{sm}}$ is an entire function and carries no singularity claim;
the finite comparisons of this section use only the restriction. The control
side therefore carries neither an Euler product nor zeros of its own: the
contrast against the primes is pure cross-coupling.
\end{remark}

\begin{theorem}[Summation formula for the smooth measure; proved, finite
ordering certified]
\label{thm:inv}
For $\varphi\in C^1[2,\kappa]$,
\[
\sum_{p\le\kappa}\varphi(p)-\int\varphi\,d\mu_{\mathrm{sm}}
=(16-M)\,\varphi(\kappa)-\int_2^{\kappa}E_\pi(t)\,\varphi'(t)\,dt ,
\]
with, at the reference cutoff, $|16-M|=2.184690$ and
\[
\sup_{t\le\kappa}|E_\pi(t)|=M-15=3.1846895084\ldots\approx3.184690,
\]
the supremum being attained as $t\uparrow\kappa$, that is immediately
\emph{before} the jump at the cutoff prime, where $E_\pi$ equals $15-M$; at
$t=\kappa$ itself it has already returned to $16-M$. The identity and the
reduction to finitely many candidates are proved; the identification of the
maximiser among them is certified in ball arithmetic by
\texttt{certify\_smooth\_controls.py} (v2.0), which encloses every candidate
and verifies the chain $\operatorname{lower}|E_\pi(c_j)|
>\operatorname{upper}|E_\pi(c_{j+1})|$ along the whole ranking, so that the
complete finite ordering of the candidates --- not only its maximiser --- is
certified.
\end{theorem}

\begin{proof}
Both sides are integrals against $d(\#\{p\le t\}-\mu_{\mathrm{sm}})=dE_\pi$;
integrating by parts on $[2,\kappa]$ and using $E_\pi(2^-)=0$ gives the
identity. The two constants are evaluations of $M$ and of $E_\pi$; for the second,
$E_\pi$ has unit upward jumps at the primes and decreases continuously in
between, so on $[2,\kappa]$ the extrema of $|E_\pi|$ occur among exactly
$31$ one-sided values: the sixteen right values $|E_\pi(p^+)|$ and the fifteen
left limits $|E_\pi(p^-)|$, $p=3,\dots,\kappa$ --- monotonicity alone decides
no more than that. Each candidate is an evaluation of the same
$\mathrm{li}$-difference and is enclosed as a ball; the directed separation
$\operatorname{lower}|E_\pi(\kappa^-)|>\operatorname{upper}|E_\pi(c)|$ then
holds along the whole ranking, so the finite ordering is certified as a total
order: the maximiser is $\kappa^-$ with margin $\ge0.1201$ over the second
candidate $41^-$, which in turn exceeds the third, $37^-$, by $\ge0.0920$; the
right endpoint contributes only $|16-M|=2.184690$.
\end{proof}

\begin{theorem}[Quadrature estimate for the discrete world; proved]
\label{thm:quad}
The position-side curvature kernel is
\[
K_\varepsilon(t):=(\log t)^2\sum_{k=1}^{N}e^{-\varepsilon^2\gamma_k^2}
\gamma_k^2\cos(\gamma_k\log t),\qquad t\in[2,\kappa],
\qquad B[\nu]:=\int K_\varepsilon\,d\nu,
\]
so that $B[P]=\sum_{p\le\kappa}K_\varepsilon(p)=B$ and

\[
\Bsm=\int_2^{\kappa}K_\varepsilon(t)\,\frac{dt}{\log t}.
\]
 With this
kernel, so that
$B[\nu]=\int K_\varepsilon\,d\nu$ for any measure $\nu$ on $[2,\kappa]$, and
$y(t)=\int_2^t ds/\log s$ the logarithmic-integral coordinate,
\[
B[A]=\frac{16}{M}\,\Bsm+E,
\qquad
|E|\;\le\;\frac{16}{M}\cdot\frac{M^3}{6144}\,
\sup_{y\in[0,M]}\Bigl|\frac{d^2}{dy^2}(K_\varepsilon\circ t)(y)\Bigr| ,
\]
\[
\frac{16}{M}\cdot\frac{M^3}{6144}
=\frac{M^2}{384}
\le 0.861154 ,
\]
the mass wedge $M/16=1.136543$ being an exact normalisation factor and not an
error term. Equivalently, before the mass rescaling,

\[
|\Bsm-(M/16)B[A]|\le(M^3/6144)\sup|d^2(K_\varepsilon\circ t)/dy^2|
\le 0.978739\,\sup|\cdot|
\]
: the two constants are one bound in two
normalisations, and the derivative is taken in $y$, the variable the rule
integrates, not in $t$.
\end{theorem}

\begin{proof}
Substitute the logarithmic-integral coordinate
$y(t)=\int_2^t ds/\log s$, which maps $[2,\kappa]$ onto $[0,M]$ and carries
$\mu_{\mathrm{sm}}$ to Lebesgue measure. In that coordinate the nodes are the
midpoints $y_j=(j-\tfrac12)M/16$ of $16$ equal panels, so $A$ is exactly the
midpoint rule for $\int_0^M K_\varepsilon(t(y))\,dy$. The midpoint rule on $n$
panels of total length $L$ has error at most
$L^3/(24n^2)\sup|\partial_y^2(K_\varepsilon\circ t)|$, which with $L=M$ and
$n=16$ is $M^3/(24\cdot256)$ times that supremum; multiplying by the mass factor
$16/M$ gives the stated constant. The second derivative is the one taken in
$y$, not in $t$: it is the composed kernel that the rule integrates.
\end{proof}

The bound of Theorem~\ref{thm:quad} is not sharp for oscillatory integrands.
It assumes worst-case phases and uses no cancellation, which for the curvature
kernel costs about five orders of magnitude at small $\varepsilon$: sixteen
nodes cannot resolve frequencies $\gamma\gtrsim 1/\varepsilon$. This is
declared, not hidden, and it is the reason the discrimination below is
certified directly rather than deduced from the quadrature estimate.

\begin{theorem}[Certified discrimination in curvature; certified]
\label{thm:disc}
Uniformly for $\varepsilon\in[0.04,0.07]$,
\[
|B[P](\varepsilon)-\Bsm(\varepsilon)|\;\ge\;5343.90,
\qquad
|B[P](\varepsilon)-B[A](\varepsilon)|\;\ge\;5416.28 .
\]
At the reference width the certified ball values are
$B[P](0.05)=-19342.5476062333843$, $\Bsm(0.05)=-241.150637407723$ and
$B[A](0.05)=-458.627462650173$, so at the reference width the smooth measure reproduces
$\approx1.247\%$ of the prime curvature. We write $\Ddisc=5416.28$ for the margin of the discrete comparison, which is
the one entering Theorem~\ref{thm:fact}; the comparison against
$\mu_{\mathrm{sm}}$ has the smaller margin $5343.90$ and is quoted for
orientation only.
\end{theorem}

\emph{Certification.} The quoted margins are \emph{downward-rounded
certified lower bounds}: ball arithmetic in Arb at $192$ bits of
precision~\cite{Johansson17}, with the $\varepsilon$-dependence separated
exactly, the smooth-measure integrals evaluated through a closed
antiderivative rather than quadrature, a covering of the window by $800$
subintervals with exact rational endpoints, and explicit bounds on the
$\varepsilon$-derivative. The ordinate inputs are certified by Arb
zeta-zero isolation under exact double-precision cross-check against the
normative layer, with $100$ Hardy-$Z$ sign-change enclosures of radius
$10^{-20}$ retained as a redundant second check.
The artefact is named in \S\,\ref{sec:methods}.

\begin{theorem}[Factorisation barrier: bounded-Lipschitz resolution lower
bound; proved from certified inputs]
\label{thm:fact}
Let
\[
\mu_P=\sum_{p\le\kappa}(\log p)^2\delta_p,
\qquad
\mu_A=\sum_{j\le16}(\log a_j)^2\delta_{a_j}
\]
be the $(\log)^2$-weighted
counting measures on $[2,\kappa]$, the weighting being the one the curvature
itself carries; their masses are $141.6150$ and $139.0471$. Write

\[
\hat K_\varepsilon(t):=K_\varepsilon(t)/(\log t)^2
\]
 for the reduced kernel,
so that $B[P]=\int\hat K_\varepsilon\,d\mu_P$ and
$B[A]=\int\hat K_\varepsilon\,d\mu_A$: the $(\log)^2$ weight sits in the
measures here and in the kernel of Theorem~\ref{thm:quad}, never in both. A
factoring functional is a map $\Psi$ on finite signed measures on
$[2,\kappa]$; equip that space with the flat bounded-Lipschitz distance
\[
\dBL(\mu,\nu):=\sup\Bigl\{\int\varphi\,d(\mu-\nu)\ :\
\|\varphi\|_\infty\le 1\ \text{ and }\ \mathrm{Lip}(\varphi)\le 1\Bigr\},
\]
the flat convention, under which the Dudley sum variant
$\|\varphi\|_\infty+\mathrm{Lip}(\varphi)\le1$ gives a smaller value, so the
conclusion holds a fortiori. Then, with $\dBL(\mu_P,\mu_A)$ as certified in
Lemma~\ref{lem:dbl} below, no functional of the counting data with
bounded-Lipschitz constant
\[
L\;<\;\frac{\Ddisc}{\dBL(\mu_P,\mu_A)}
\]
can factor the discrimination of Theorem~\ref{thm:disc}.
\end{theorem}

For this single pair of measures the bound is exactly the pairwise
Lipschitz quotient; no stronger obstruction to factorisation is claimed.

The distance itself is a certified input, stated separately so that the
theorem's status is exact.

\begin{lemma}[Bounded-Lipschitz distance of the weighted measures;
certified]
\label{lem:dbl}
$\dBL(\mu_P,\mu_A)=152.529455\ldots$, rigorously two-sided.
\end{lemma}

\emph{Certification.} The value is obtained as a linear program in the
Kantorovich--Rubinstein form, solved to optimality and certified in two
directions: primal feasibility of the returned $\varphi$ is verified in ball
arithmetic, and weak duality with residual absorption
($\|r\|_1=1.8\cdot10^{-14}$) bounds the optimum from the other side.

\begin{proof}[Proof of Theorem~\ref{thm:fact}]
If $B[P]=\Psi(\mu_P)$ and $B[A]=\Psi(\mu_A)$ for a functional $\Psi$ with
bounded-Lipschitz constant $L$, then

\[
|B[P]-B[A]|=|\Psi(\mu_P)-\Psi(\mu_A)|\le L\,\dBL(\mu_P,\mu_A).
\]

Combining with the lower bound of Theorem~\ref{thm:disc} and the certified
value of Lemma~\ref{lem:dbl} gives $L\ge\Ddisc/\dBL$.
\end{proof}

The threshold depends on which margin is inserted, and the three figures in
circulation are all the same statement at different strengths. With the margin
declared in advance of the computation, $5000$, the threshold is
$5000/152.529455=32.78$; with the margin actually achieved uniformly on the
window, $\Ddisc=5416.28$, it is $35.51$; and pointwise, using the realised
difference at a single width instead of the uniform bound, it is $123.8$ at
$\varepsilon=0.05$ and $283.8$ at $\varepsilon=0.04$. The sharp uniform value
$35.51$ is the one quoted in the abstract.

The honest reading of Theorem~\ref{thm:fact} is a resolution lower bound, not
an impossibility statement. The curvature is itself bounded-Lipschitz, with
kernel sup-norm at most $664$ at the reference width; the theorem says that
any factoring functional must have a constant at least in the tens, that is,
must resolve the counting data far more finely than the density does. The
statement covers one pair of worlds; the family form is open, and functionals
outside the bounded-Lipschitz class are untouched.

The programmatic consequence is the one that matters for
\S\,\ref{sec:audibility}, and two levels must be kept apart before it is
drawn: the ordinate channel is external and identical in both comparisons,
while the intrinsic zeta function of the control side is a separate object ---
it is the former that carries the coupling. On that external channel both
objects compared here have the same zero side, and the intrinsic zeta
function of the smooth measure has no zeros at all, so the entire
discrimination is carried by the prime side. What is certified is precise and
should be stated precisely: the curvature separates from the smooth measure
with uniform margin $\ge5343.90$ and from the discrete world with uniform
margin $\ge5416.28$, and for the discrete pair --- the pair the
bounded-Lipschitz distance is attached to --- no functional of the counting
data with constant below the stated threshold factors that separation. That
identifies the prime-side weighted point configuration, relative to the
stated controls, as audible; it does not by itself identify multiplicativity
as the cause, because the controls vary the point geometry and the
multiplicative structure together.

\section{Limits of Audibility}
\label{sec:audibility}
% ============================================================

The barrier comes first.

For the curvature sums the relevant carrier is the aggregated curvature
test kernel
\[
\mathcal K_{\kappa,\varepsilon}(z)
:=\sum_{p\le\kappa}(\log p)^2\,z^2 e^{-\varepsilon^2z^2}\cos(z\log p),
\]
a sum over \emph{primes} in the \emph{ordinate} variable --- distinct from
$G$, which sums over ordinates in the ratio variable; the two must not be
interchanged. That $\mathcal K_{\kappa,\varepsilon}$, and not $G$, carries
the continued curvature sums is an identity, not a reading: at the reference
data,
\[
\sum_{k=1}^{N}\mathcal K_{\kappa,\varepsilon}(\gamma_k)=B=-19342.5476,
\]
the per-ordinate decomposition of the curvature anchor of the series.

\begin{theorem}[Functional-equation symmetry barrier; proved]
\label{thm:barrier}
Let a zero be written $\rho=\beta+i\gamma$ and let $\delta=\beta-\tfrac12$
denote its distance from the critical line. The functional equation
$\xi(s)=\xi(1-s)$, together with reality of $\xi$ on the real axis, places
zeros in orbits $\{\tfrac12\pm\delta\pm i\gamma\}$. Every symmetric functional
of such an orbit is an even function of $\delta$. Evenness holds without any
regularity hypothesis, and it already implies that no symmetric functional of
the zero data --- in particular none of the observables of this paper, and
none of the curvature sums of the earlier papers --- can determine the
\emph{sign} of $\delta$. If such a functional is in addition differentiable at
$\delta=0$, its first derivative there vanishes; if it is analytic, the first
possibly non-zero Taylor coefficient is of even order. For the curvature sums
of this series, whose $\delta$-dependence enters through the even orbit
average

\[
\tfrac12\bigl(\mathcal K_{\kappa,\varepsilon}(\gamma_k+i\delta)
+\mathcal K_{\kappa,\varepsilon}(\gamma_k-i\delta)\bigr),
\]

that first term is the quadratic one, unless the aggregate quadratic
coefficient happens to cancel. The class of observables studied here is
therefore blind to the orientation of the displacement ---
unconditionally. When such an observable is differentiable in $\delta$
at $0$, its first derivative there vanishes; when it is real-analytic, its
expansion contains only even powers, and the response to the magnitude
begins at even order.
\end{theorem}

\begin{proof}
The orbit is invariant under $\delta\mapsto-\delta$ by the functional
equation combined with conjugation. A symmetric functional of an
invariant set is invariant under any permutation of it, hence takes the same
value at $\delta$ and $-\delta$; an even function of $\delta$ has vanishing
first derivative at $\delta=0$. For the curvature sums the mechanism is
explicit: each ordinate enters the continued sums through the even orbit
average
\[
\tfrac12\bigl(\mathcal K_{\kappa,\varepsilon}(\gamma_k+i\delta)
+\mathcal K_{\kappa,\varepsilon}(\gamma_k-i\delta)\bigr)
=\mathcal K_{\kappa,\varepsilon}(\gamma_k)
-\tfrac{\delta^2}{2}\,\mathcal K_{\kappa,\varepsilon}''(\gamma_k)
+O(\delta^4),
\]
so the first-order term vanishes identically and the first possibly non-zero
contribution is the aggregate quadratic one.
\end{proof}

This is a limit, not a lament, and it is stated before the measurement so that
the measurement cannot be misread. The statement is precise about what fails:
not that the observables carry no trace of a displacement whatever, but that
they cannot resolve its direction, and that, under differentiability, the
first-order response to its magnitude vanishes --- for the real-analytic
curvature sums of this series, the response begins at the quadratic
term, unless the aggregate quadratic coefficient happens to cancel. What
the theorem bars is any inference to the
\emph{orientation} of the displacement. The even-order response can depend on
$|\beta-\tfrac12|$, so an inference to the \emph{magnitude} is not excluded by
the theorem; whether the magnitude is in fact reconstructible from these
observables is neither proved nor tested in this paper, and no such inference
is made anywhere in it.

What remains open to measurement is the ordinate side: does the coupling hear
anything about the ordinates beyond their density class? The following
measurement was designed and registered before it was run.

\emph{Normalisation.} The ordinate axis is varied while the prime side is held
fixed at $p\le\kappa$. Ordinates are unfolded against each world's own smooth
counting function, $\tilde b_k:=N_W(b_k)$, giving unit density, and the
Gaussian weight is applied on the unfolded scale with
$\tilde\varepsilon=0.25$; robustness is checked at $\tilde\varepsilon=0.10$.
Only unfolded statements carry discriminating content beyond density. The
phasor argument in this battery is $\tilde b\log p$, without the trace
parameter; this differs from the convention of \S\S\,\ref{sec:operator} to
\ref{sec:barrier} and is declared here because the observables below are
sensitive to it. The convention was fixed before the computation, not chosen
after it.

\emph{Design.} Five worlds, each constructed independently and given
explicitly: the Riemann ordinates $W_\zeta$; the ordinates $W_{\chi_4}$ of the
Dirichlet $L$-function of the non-principal character modulo $4$; those of the
Dedekind zeta function of $\mathbb{Q}(i)$, a degree-two Euler object assembled
as a union, $W_{\mathrm{Ded}}$; those of the Davenport--Heilbronn
function~\cite{DavenportHeilbronn36}, which has a functional equation but no
Euler product and of which only the imaginary ordinates enter --- its off-line
real parts are discarded --- $W_{\mathrm{DH}}$; and a synthetic density skeleton
$W_{\mathrm{syn}}$ with eight replicates supplying the scatter scale $S_i$.
The skeleton is the exactly unfolded sequence $\tilde b_k=k-\tfrac12$,
$k=1,\dots,30$, and the replicates are
\[
\tilde b_k^{(m)}=k-\tfrac12+0.25\,\sin\bigl(2\pi\{k\phi_m\}\bigr),
\qquad \phi_m=\{\phi^{m}\},\quad m=1,\dots,8,
\]
with $\phi$ the golden ratio and $\{\cdot\}$ the fractional part: an explicit
deterministic recipe, with no random number generator anywhere. The recipe
degenerates, and the degeneration is worth stating precisely. Since
$\phi^2=\phi+1$, one has $\{\phi\}=\{\phi^2\}$, so the replicates $m=1$ and
$m=2$ coincide and only seven are distinct. Because $S_i$ below is a maximum
over the replicates, a repeated replicate contributes a value already present:
$S_i$ is numerically \emph{unchanged} by the degeneration. What is lost is one
degree of freedom in the design, not scatter.

With
\[
\Phi_{ij}=e^{-\tilde\varepsilon^2\tilde b_i^2/2}\sin(\tilde b_i\log p_j),
\]
$G^{\mathrm{un}}=\Phi^{\!\top}\Phi$ on the $16$-dimensional prime space and
$\Tt=\Phi\Phi^{\!\top}$ on the ordinate space, the five observables fixed in
advance were

\begin{center}
\begin{tabular}{@{}cll@{}}
\toprule
 & Observable & Normalisation \\
\midrule
$\mathcal{O}_1$ & $\sum_j(\log p_j)^2\sum_i e^{-\tilde\varepsilon^2\tilde b_i^2}\tilde b_i^2\cos(\tilde b_i\log p_j)$
   & raw \\
$\mathcal{O}_2$ & $\lambda_{\max}(G^{\mathrm{un}})/\mathrm{Tr}\,G^{\mathrm{un}}$ & dimensionless \\
$\mathcal{O}_3$ & $\mathrm{Tr}\,G^{\mathrm{un}}$ & raw \\
$\mathcal{O}_4$ & $\bigl(\sum_{i,j}|\Tt_{ij}|-\sum_i|\Tt_{ii}|\bigr)/ \sum_{i,j}|\Tt_{ij}|$ & fraction \\
$\mathcal{O}_5$ & $\lambda_2(G^{\mathrm{un}})/\lambda_1(G^{\mathrm{un}})$ & dimensionless \\
\bottomrule
\end{tabular}
\end{center}

\emph{Criterion.} Registered in advance, with
\[
D_i(W):=\bigl|\mathcal{O}_i(W)-\mathcal{O}_i(W_{\mathrm{syn}})\bigr|,
\qquad
S_i:=\max_{m=1,\dots,8}
\bigl|\mathcal{O}_i(W^{(m)})-\mathcal{O}_i(W_{\mathrm{syn}})\bigr|,
\]
a world $W$ separates on $\mathcal{O}_i$ if $D_i(W)>3S_i$. The design note fixing worlds, observables, criterion and window
is dated and archived with the verification code~\cite{Code}.

\emph{Outcome (numerical).} The criterion is not met, and not narrowly. At
$\tilde\varepsilon=0.25$ the scatter scales are
$S_1=640.995$, $S_2=0.151280$, $S_3=4.71715$, $S_4=0.0515498$,
$S_5=0.483144$, and the ratios $D_i/S_i$ are

\begin{center}
\begin{tabular}{@{}lccccc@{}}
\toprule
 & $\mathcal{O}_1$ & $\mathcal{O}_2$ & $\mathcal{O}_3$ & $\mathcal{O}_4$
 & $\mathcal{O}_5$ \\
\midrule
$W_\zeta$            & 0.299 & 0.069 & 0.055 & 0.580 & 0.135 \\
$W_{\chi_4}$         & 0.342 & 0.031 & 0.064 & 0.324 & 0.073 \\
$W_{\mathrm{Ded}}$   & 0.154 & 0.630 & 0.586 & 0.270 & 0.270 \\
$W_{\mathrm{DH}}$    & 0.082 & 0.060 & 0.018 & 0.168 & 0.012 \\
\bottomrule
\end{tabular}
\end{center}

The largest entry of the whole matrix is $0.630$, a factor $4.8$ below the
threshold; not one cell reaches even one scatter unit. At
$\tilde\varepsilon=0.10$ the picture is unchanged, with maximum $0.799$. In
particular the world with a functional equation but no Euler product does not
separate either: there is no non-Euler anomaly on this axis. The outcome is
qualitatively compatible with broad universality heuristics for unfolded local
statistics~\cite{RudnickSarnak96}, but the battery does not test an $n$-level
correlation law; its registered conclusion is only that none of the five
observables separates the declared worlds at the stated threshold.

\emph{Scope.} The finding holds for this battery, this window and this
coupling. Pair and correlation statistics of the unfolded ordinates were not
examined; any continuation would be a new experiment requiring its own
pre-registration. The effective sample behind the gate is small ---
$n=(\sum_kw_k)^2/\sum_kw_k^2\approx7.1$ at $\tilde\varepsilon=0.25$ and
$17.6$ at $\tilde\varepsilon=0.10$, the Kish effective size of the Gaussian
gate weights $w_k=e^{-\tilde\varepsilon^2\tilde b_k^2/2}$ on the $30$
unfolded ordinates (\texttt{ess\_gate.py}) --- and this was declared in
advance as the weakest point. No no-go theorem is derived from
this measurement, and it says nothing whatever about the Riemann Hypothesis.
The term \emph{pre-registered} is used here as a methodological import from
the experimental sciences, and is meant literally: design, observables and
criterion were fixed in a dated document before the computation was run.

\emph{Completeness.} That the ordinate lists are complete in the window was
verified by rectangle counting via the argument principle. For the two worlds
whose lists had to be constructed by an independent sign search --- the
character world and the world with a functional equation but no Euler product
--- the winding number came out at exactly $30$ in each case, with contour
minimum modulus at least $0.27$. This is a numerical verification, not an
interval-arithmetic one, and is labelled accordingly. As a by-product, all
zeros of the Davenport--Heilbronn function in the window lie on the line, the
first off-line ordinate, at $85.70$~\cite{Spira94}, lying above the window
maximum $64.04$ --- a statement about the known zero list, independent of the
rectangle count.

\section{The Audibility Classification and Discussion}
\label{sec:discussion}
% ============================================================

The three preceding sections classify the coupling along three axes.
Throughout this paper, \emph{audible} is used in a relative sense: a feature
of the data is called audible if some observable of the stated class
separates the world carrying it from the stated comparison worlds. The
classification below is therefore always relative to the declared observable
class and the declared controls; no absolute notion of audibility is
asserted.

\begin{center}
\begin{tabular}{@{}l p{4.6cm} l@{}}
\toprule
Axis & Finding & Status \\
\midrule
Prime-side point configuration vs.\ stated controls & audible, margins
  $\ge5343.90$ / $\ge5416.28$ & certified \\
Prime-side Beurling within-class variation & \raggedright not distinguished at the
  level of the proved leading laws; same $R$, $\Delta_0$, $F$ and band;
  lower-order discrimination not excluded & proved \tabularnewline
Ordinates beyond unfolded one-point density & not separated on this battery &
numerical \\
Zero-orbit orientation & not audible; first derivative zero when
  differentiable & proved \\
\bottomrule
\end{tabular}
\end{center}

The order is the point. The barrier is a theorem, so the last row is not a
gap in the measurement but a boundary of the object; the third row is a
measurement with a declared scope; and the second row is a proved
\emph{invariance} --- within the class the proved leading asymptotic data do not distinguish:
$R$, $\Delta_0$, $F$ and the band are class-independent. The
system-dependent remainder terms are not controlled sharply enough here to
exclude lower-order discrimination. The first row states what is certified ---
the prime-side point configuration separates from the stated controls --- and
a multiplicative reading of that separation, that what is heard is Euler
structure, is available as an \emph{interpretation} and is marked as one here:
the controls of \S\,\ref{sec:barrier} vary point geometry and multiplicative
structure together, so the certified statement does not attribute the
separation to either alone. The identification experiment that would decide
this is stated as Open Problem~\ref{op:ident}. A function-side
reconstruction experiment, reported after this discussion, is directionally
compatible with this reading. Read so, the first and
second rows together say that this coupling is an instrument for arithmetic
point structure and for counting density --- the latter as the surviving
class parameter --- and for nothing else that we were able to test on this
finite battery.

\emph{Numerical observation} (function-side reconstruction; a separate
experiment, not part of the \S\,\ref{sec:audibility} battery). The
experiment asks how much of the zero placement, and of the area signature
between consecutive zeros, a phasor field recovers along the critical line.
Truth is Hardy's $Z(t)$ on $t\in[100,600]$ (grid $0.04$; $312$ zeros of the
extended list \texttt{zeros\_650.csv}, cross-checked against the normative
first-$100$ layer, which it reproduces exactly in double precision, hence
$311$ areas); the models are the finite cosine sums
$F(t)=\sum_n 2n^{-1/2}\cos(\theta(t)-t\log n)$ over declared index sets
inside the Riemann--Siegel window $n\le\sqrt{t/2\pi}$, with two
pre-declared metrics: zero recovery at tolerance $0.25$ of the local mean
spacing, and Spearman rank correlation of the areas. With all amplitudes
fixed at $2n^{-1/2}$, the full window set recovers $98.1\%$ of the zeros
with $0.0\%$ spurious crossings and area correlation $0.980$; prime indices
alone drop recovery to $34.6\%$ with $35.4\%$ spurious; adding the constant
mode and the prime powers $\{4,8,9\}$ reaches $85.9\%$; the single
semiprime term $n=6=2\cdot3$ closes the gap back to $98.1\%$. After the
amplitude normalisation, the degradation pattern under removal of composite
modes is consistent with sensitivity to combination frequencies of already
active prime modes. The comparison with the curvature side is directional
only: the curvature controls of \S\,\ref{sec:barrier} vary multiplicative
structure and point geometry together, so the two experiments are
directionally compatible, but no common causal feature is established. The
experiment is produced by \texttt{function\_side\_reconstruction.py} (\S\,\ref{sec:methods}).

\subsection{Near-coincidence fusions}

\emph{Numerical observation.} The resolution scale of
\S\,\ref{sec:operator} has a sharper consequence once
the mode set is completed to prime powers. Three further pairs fuse: $31$ with $32$ at
correlation $0.977$, $16$ with $17$ at $0.908$, and $47$ with $49$ at $0.907$.
In each case a prime power lands beside a neighbouring prime --- $2^5$ beside
$31$, $2^4$ beside $17$, $7^2$ beside $47$ --- and the two directions become
indistinguishable to the operator. The effect is deterministic logarithmic
geometry and nothing else: the kernel sees $\log n$. Across the completed
$24$-mode reference set, and at the reference parameters, every pair whose
logarithms differ by less than roughly $0.065$ was found unresolved; this is
a one-sided empirical rule for that set, not a general resolution criterion,
since the pair correlation runs through $G(\log mn)$ as well and is
therefore not monotone in $|\log m-\log n|$ alone. This threshold is genuinely smaller than the rule of
thumb $0.12$ obtained by calibrating on prime pairs alone, which is too loose
for power pairs; the pair $(25,27)$ at logarithmic distance $0.0770$ does not
fuse, reaching only $0.540$, and the pair $(8,9)$ at $0.1178$ reaches
$0.224$. Multiplicative structure meets additive structure at these points,
and the instrument cannot tell them apart.

\subsection{Delineation}

The following delineations are stated rather than implied.

\emph{Prime cutoff against prime indexing.} The line of work of Connes and
collaborators constructs self-adjoint operators from truncated Euler products,
the truncation parameter controlling which primes $p\le c$ enter the
operator~\cite{Connes26,ConnesVS25,Groskin26a}. Self-adjointness is there by
construction, so the spectrum is real by design; what is approximated is the
sequence of ordinates, and the entire weight of the Riemann Hypothesis rests
on an unproven convergence as the cutoff grows. This does not conflict with
Theorem~\ref{thm:barrier}, whose subject is different: not whether some
operator can be built whose spectrum is real, but whether the distance of a
zero from the critical line can be \emph{read off} a class of observables that
presupposes nothing. In the language of that programme, the scaling spectral
triple plays the role of the archimedean factor --- the truncated completion
$\xi=\Gamma_{\mathbb{R}}\prod(1-p^{-s})^{-1}$ in operator form --- and the same
ingredient that produces the zeros removes, for symmetric functionals, their
location. Here the primes are the index set of a Gram operator and the
ordinates are frequencies; that is a different architecture, not a variant of
the same one.

\emph{The nearest neighbour.} Groskin's finite Guinand--Weil
dictionary~\cite{Groskin26b} establishes, for the truncated Weil quadratic
form, an exact finite correspondence in which the value of the truncated form
on a Galerkin coefficient vector equals a sum of an explicitly constructed
band-limited test function over the non-trivial zeros. This is thematically
adjacent to the trace identity of \S\,\ref{sec:operator}: both are exact
finite bridges between an operator quantity and a sum over zeros.
Architecturally the two are separate. That construction lives on the
test-function side, in a frequency-band basis, and its object is Weil
positivity; there is no prime--ordinate Gram coupling in it, no variation of
the arithmetic world, and no question of audibility. Neither result implies
the other, and neither is a special case of the other.

\emph{Screw functions.} Suzuki's unification of the Weil quadratic
form~\cite{Suzuki26} treats the same positivity question through continuous
functions rather than distributions; it is adjacent, and likewise free of any
assumption of the Riemann Hypothesis. It does not index by primes.

\emph{Local Euler factors.} Liu's operator pencils~\cite{Liu26} encode local
Euler factors prime by prime and compose them as a restricted tensor product.
Primes do appear there as an index family, which makes this the closest
parallel track on the indexing question; but the objects are local factors,
not a global coupling between primes and ordinates, and no spectral ratio,
no class transfer and no audibility question arises.

\emph{Arithmetical matrices.} Hilberdink and Pushnitski~\cite{HilberdinkPushnitski24}
obtain spectral asymptotics for a fixed family of arithmetical matrices, with
the same constant governing both spectral branches, and indicate a connection
to Beurling systems. That constant belongs to that matrix family. Theorem~\ref{thm:class}
establishes a different statement: the present leading objects $R$,
$\Delta_0$, $F$, and their band are invariant across systems in
$\Dclass$. It does not establish uniqueness of those objects across
different density classes.

\emph{Operators over Beurling-induced spaces.}
Kouroupis~\cite{Kouroupis23} studies Hardy spaces of generalised Dirichlet
series induced by Beurling primes and shows that boundedness and compactness
of certain composition operators do not depend on the choice of primes. This is genuine operator
theory over Beurling systems, but there is no Gram matrix, no ordinate
channel, and no fingerprint in asymptotic spectral constants.

\emph{Off-line zeros.} The Davenport--Heilbronn function is the standard
reference for the statement that a functional equation without an Euler
product does not force the zeros onto the line~\cite{DavenportHeilbronn36}. It
is not the first documented off-line zero: that record belongs to the work of
Potter and Titchmarsh on the Epstein zeta function~\cite{PotterTitchmarsh35,
Epstein03}.

% ============================================================
\section{Methods and Verification}
\label{sec:methods}

\emph{Ordinate protocol.} Two layers are used and never mixed. The normative
layer is a fixed list of the first $100$ ordinates in double precision; it
defines every anchor value and every tolerance quoted in this paper, and it is
the reason no figure here is quoted beyond the precision that list supports.
Where a certificate requires more, the ordinates are obtained as certified
enclosures from Arb's zeta-zero isolation, and the ball midpoints are subject
to a mandatory cross-check against the normative layer, which they must
reproduce exactly in double precision; the enclosures are then part of the
certificate itself. For non-certified high-precision work the regeneration
recipe remains \texttt{mpmath.zetazero} at $30$ digits, under the same
cross-check. No regenerated list is archived; the generating scripts are.
The extended list \texttt{zeros\_650.csv} is a named, non-normative
auxiliary file used for exactly one numerical experiment
(\S\,\ref{sec:discussion}); it defines no anchor and no tolerance, and its
first $100$ entries reproduce the normative layer exactly in double
precision. No further \emph{normative} ordinate file exists.

\emph{Artefacts.} Each machine-checked statement of this paper has a named
artefact in the code repository~\cite{Code}. The anchor and prohibition checks
of the manuscript are \texttt{verify\_paper9.py} (v0.30, 277 checks;
the script verifies its own check count and exits with failure on any
shortfall, so a partial run cannot end in success --- it is a drift gate over
the quoted values, recomputing the anchors marked for recomputation, and does
not itself reproduce the experiments). The five-world battery of
\S\,\ref{sec:audibility} is produced by \texttt{unfolded\_discrimination.py}:
world construction including the Hurwitz-zeta $\chi_4$-line and the
Davenport--Heilbronn ordinates after the functional-equation construction,
the five observables with their scatter scales, and the completeness counts
by the argument principle; the non-$\zeta$ ordinate lists are constructed
deterministically inside the script, which therefore serves as their
generation protocol; for $W_\zeta$ the script regenerates the first $31$
ordinates internally with \texttt{mpmath.zetazero} at $40$ digits and
cross-checks them against the normative \texttt{zeros\_100.csv}. The effective gate
size is the one-line computation of \texttt{ess\_gate.py}. Both
certificate artefacts obtain their ordinates by Arb zeta-zero isolation
(\texttt{acb.zeta\_zeros}) under exact double-precision cross-check
against the normative layer. The function-side reconstruction experiment of \S\,\ref{sec:discussion} is
produced by \texttt{function\_side\_reconstruction.py}: truth grid, index sets, tolerance and
both metrics are declared in the script, and the quoted recovery and
correlation figures are its direct output; its ordinate input is the
auxiliary file \texttt{zeros\_650.csv}. The two-point separation of
Theorem~\ref{thm:nonconv} is certified by \texttt{cert\_paper9\_ratio.py}
(v1.2), which obtains the ordinates as certified enclosures from Arb's
zeta-zero isolation, cross-checks the ball midpoints against the normative
double-precision layer, builds the limit matrices as ball matrices, bounds the
leading eigenvalue from below by a Rayleigh quotient at $u_1$ and from above
by a trace power at $u_2$, and divides the resulting intervals outward. The margins of Theorem~\ref{thm:disc}, the closed
antiderivative replacing quadrature, the window covering and the linear
program certifying $\dBL$ in Theorem~\ref{thm:fact} are produced by
\texttt{certify\_smooth\_controls.py} (v2.0) as downward-rounded certified lower bounds;
for non-certified high-precision regeneration,
\texttt{gen\_zeros.py} computes the first $100$ ordinates with
\texttt{mpmath.zetazero} at $30$ digits and cross-checks their
double-precision values against the normative \texttt{zeros\_100.csv}. It
is not an input to either Arb certificate. Each certificate script obtains
its own certified zero enclosures internally via \texttt{acb.zeta\_zeros}
and performs the mandatory midpoint cross-check against the normative
layer.

\emph{Certificate stack.} Certified statements use ball arithmetic in Arb
through its Python interface, \texttt{python-flint}~0.9.0 on
FLINT~3.6.0~\cite{Johansson17}; non-certified numerics use
\texttt{mpmath}~1.4.0~\cite{mpmath} at the precision stated in each place. The
off-line control quoted from the earlier work of this series rests on the
verified height $H_0=3\cdot10^{12}$~\cite{PlattTrudgian21}, an unconditional
finite computation.

\emph{Reproducibility.} The mechanical checks are of four kinds, and the
gate distinguishes them: values marked for recomputation are recomputed from
the ordinate list; textual anchors are drift anchors --- their presence is
enforced, their values fixed by the artefact that produced them; interval
statements have dedicated certificate artefacts whose outputs are themselves
checked as anchors; and the identities of Theorems~\ref{thm:imported}
and~\ref{thm:kernel} are checked as identities, not as approximate
agreements. Plots are not reproduced here; they live in the
code repository together with the scripts.

\section{Open Problems}
\label{sec:open}
% ============================================================

These are bequests, not a programme. Each is stated in the terms in which it
would have to be answered.

\begin{openproblem}[Pointwise finality of the divergence]
\label{op:pointwise}
Theorem~\ref{thm:eta} is a statement about window means, and
Theorem~\ref{thm:nonconv} about limit points.

\smallskip
\noindent\emph{Given.} The profile $f(e^u)$ of Lemma~\ref{lem:fmean} and the
recurrence of the Kronecker flow.

\smallskip
\noindent\emph{Target.} A pointwise statement: does
$\eta_{\mathrm{orig}}(\kappa)\to-\infty$ without averaging?

\smallskip
\noindent The question splits into two tasks. First, remainder control in
the small-profile regime: pointwise, Step~3 of the proof of
Theorem~\ref{thm:eta} gives
\[
\eta_{\mathrm{orig}}(e^u)=-\frac{4e^uf(e^u)}{\WeN u^2}+O\Bigl(\frac{e^u}{u^3}\Bigr),
\]
so once
$f(e^u)\lesssim1/u$ the leading term falls to the size of the remainder and
the remainder decides the sign. Second, quantitative
simultaneous approximation: under the additional, unproved hypothesis that
the reference ordinates are linearly independent over $\mathbb{Q}$,
Kronecker approximation gives $\inf_u f(e^u)=0$, but supplies no rate;
without that hypothesis even the infimum is not established here. What
matters is therefore how often and how deeply $f$ nearly vanishes, not that
it does. The dichotomy is unresolved in both
directions.
\end{openproblem}

\begin{openproblem}[Minimal order break]
\label{op:order}
Theorem~\ref{thm:class} covers density classes whose remainder is
$O(x/(\log x)^{2+\delta})$. Determine the minimal deviation from
$\vartheta_{\mathfrak{B}}(x)\sim x$ that destroys the class laws, rather than merely
degrading their rates. Whether the scaling laws hold uniformly over the
sideband family of Theorem~\ref{thm:sideband} --- non-constancy of the
transferred ratio and positivity of the transferred profile for all
admissible $(a,\omega)$ --- is likewise open. So is the status of the
early-data bound: whether $M_0$ in the uniformity clause of
Theorem~\ref{thm:class} can be dispensed with at fixed $(C,x_0)$: whether
the uniformity of the conclusions persists on $\Dclass(C,x_0)$ without the
early-count bound is open; any proof would have to bypass uniform
equidistribution control, which finite insertion near $1$ destroys (see
the Remark after Theorem~\ref{thm:class}). It is not
decided here.
\end{openproblem}

\begin{openproblem}[Exceptional measure under intermediate conditions]
\label{op:exceptional}
It is not established here that the window bound admits any exceptional-set
control under the hypothesis of Definition~\ref{def:beurling}; the tail
available to the present argument is polynomial rather than summable, which
would cap such a control, if one exists, at density tending to zero. Determine
whether an exceptional-set statement holds at all under $\Dclass$, of what
strength, and what additional summability of the error term, if any, yields
finite measure.
\end{openproblem}

\begin{openproblem}[Family form of the factorisation barrier]
\label{op:family}
Theorem~\ref{thm:fact} compares one pair of worlds. Establish a barrier
uniform over a family of control worlds, and determine whether the threshold
constant can be pushed beyond the bounded-Lipschitz class.
\end{openproblem}

\begin{openproblem}[Pair statistics on the unfolded axis]
\label{op:pair}
The measurement of \S\,\ref{sec:audibility} used one-point observables.
Whether pair or higher correlation statistics of the unfolded ordinates
separate arithmetic worlds from a density skeleton, under a criterion fixed in
advance, is untested here.
\end{openproblem}

\begin{openproblem}[Identification design for the prime-side separation]
\label{op:ident}
The certified discrimination of \S\,\ref{sec:barrier} separates the prime
configuration from density-matched controls; it does not identify which
feature carries the separation, because those controls vary the one-point
geometry and the multiplicative structure together. Construct a family of
control worlds matched to the primes in multiplicative structure while varying
the one-point geometry, or matched in one-point geometry while varying
multiplicativity, with margins and criterion fixed in advance --- the
experiment that would turn the interpretation of
\S\,\ref{sec:discussion} into a result, in either direction.
\end{openproblem}

\begin{openproblem}[Canonicity beyond the trigonometric class]
\label{op:canon}
Lemma~\ref{lem:canonicity} is proved for finite spectra. The analytic
ingredients \eqref{eq:addAC} and \eqref{eq:multAC}, and the exclusion
$0\notin\mathcal{F}$, hold in the full Bohr class by
Proposition~\ref{prop:resonance} under (H1)--(H2). What is missing is purely combinatorial: do
these three conditions force $\#\mathcal{F}=1$ for countable $\mathcal{F}$
with square-summable coefficients? The hardest model case is
$\mathcal{F}\subset\mathbb{Z}$, where the conditions become fibre relations
$\sum_m a_{pm}\overline{a_{qm}}=0$. Proposition~\ref{prop:resonance}(b),
under (H1)--(H2), already excludes every spectrum in which some difference or
some same-sign ratio is uniquely represented.
\end{openproblem}

% ============================================================
\section*{Non-Circularity}
\label{sec:nocirc}
% ============================================================

We document explicitly that the results of this paper do not rely on the
objects they are directed toward.

\textbf{No RH as input.}
The Riemann Hypothesis is never assumed. The ordinates enter as numerically
specified inputs; no property of their distribution requiring the Riemann
Hypothesis is used. Theorem~\ref{thm:barrier} is a statement about the
functional equation alone.

\textbf{No GUE, Montgomery, or Hilbert--P\'olya hypothesis.}
The distribution of the ordinates is not modelled by random matrix theory.
The random-matrix literature enters only as background motivation for the
comparison reported in \S\,\ref{sec:audibility}, never as an assumption and
never as a hypothesis the battery tests. The coupling operator is neither postulated nor used as a
Hilbert--P\'olya operator; Paper~4~\cite{Paper4} reports a numerical
diagnostic failure on its tested finite-cutoff grid. Its entries are
explicit deterministic arithmetic expressions.

\textbf{No inference from a discriminator to a location.}
The discrimination established in \S\,\ref{sec:barrier} concerns the prime
side, and two levels must be kept apart. The ordinate channel that carries the
coupling is external and identical in both comparisons; separately, the
intrinsic zeta function of the smooth control world is zero-free
(Remark~\ref{rem:smoothzeta}), which is a property of that auxiliary object
and is never transferred to the external ordinates. No statement about the
position of any zero is drawn from any observable in this paper. Theorem~\ref{thm:barrier} states in which sense such
an inference is barred: the observables are even in the signed displacement, so
its orientation is not recoverable and its magnitude is invisible to first
order.

\textbf{Synthetic worlds are comparison data.}
The constructed worlds of \S\,\ref{sec:audibility} are comparison data. No
property of any of them is transferred to the zeta function, and the
Davenport--Heilbronn data are used as ordinates only.

\textbf{No claim of an RH criterion.}
Nothing in this paper is offered as a criterion for, or as evidence for or
against, the Riemann Hypothesis. The negativity of the curvature is a
one-sided second-moment fact; the classification is a statement about the
reach of an instrument.

% ============================================================
\section*{Acknowledgments}
% ============================================================

The author thanks the readers who provided feedback on earlier versions of
this paper and on the preceding papers in the series.

% ============================================================
\section*{Call for Feedback}
% ============================================================

This paper is part of an open research initiative on prime--zero
coupling and explicit-formula operators.
All papers in the series are freely available on Zenodo under CC~BY~4.0, and
the accompanying verification code is hosted on GitHub.

{\sloppy
Zenodo (\url{https://zenodo.org/search?q=Ulrich+Tehrani}) and GitHub
(\url{https://github.com/utehrani/analysislab-nt}).
\par}

Topics of particular interest include: the combinatorial question of
Open Problem~\ref{op:canon}; sharpenings of the factorisation barrier beyond
the bounded-Lipschitz class; and any observable of the coupling data that
separates an arithmetic world from a density skeleton under a criterion fixed
in advance.

% ============================================================
% References
% ============================================================
\newpage
\begin{thebibliography}{99}
\sloppy

\bibitem{Paper1}
U.~Tehrani,
``A Curvature Decomposition of the Explicit Formula
for the Riemann Zeta Function'',
Zenodo, \url{https://doi.org/10.5281/zenodo.19025598}, 2026.

\bibitem{Paper2}
U.~Tehrani,
``From Local Curvature to the Weil Functional:
An Explicit Construction'',
Zenodo, \url{https://doi.org/10.5281/zenodo.19106992}, 2026.

\bibitem{Paper3}
U.~Tehrani,
``A Finite-Cutoff Hilbert-Space Model
for Prime--Zero Energy Structure'',
Zenodo, \url{https://doi.org/10.5281/zenodo.19307989}, 2026.

\bibitem{Paper4}
U.~Tehrani,
``A Dual Operator for Prime--Zero Coupling
and a Conditional Proof of Energy Asymmetry'',
Zenodo, \url{https://doi.org/10.5281/zenodo.19364703}, 2026.

\bibitem{Paper5}
U.~Tehrani,
``Spectral Trace Formula and Smoothed Zero Sums:
A Prime--Zero Duality Framework'',
Zenodo, \url{https://doi.org/10.5281/zenodo.19508547}, 2026.

\bibitem{Paper6}
U.~Tehrani,
``Positive Curvature of the Spectral Trace
at the Critical Line'',
Zenodo, \url{https://doi.org/10.5281/zenodo.19665790}, 2026.

\bibitem{Paper7}
U.~Tehrani,
``Conditional Stationarity and Positive Curvature of the
Spectral Trace at the Critical Line'',
Zenodo, \url{https://doi.org/10.5281/zenodo.20440671}, 2026.

\bibitem{Paper8}
U.~Tehrani,
``Unconditional Negativity of the Second-Moment Bias
and Off-Line Robustness'',
Zenodo, \url{https://doi.org/10.5281/zenodo.20792123}, 2026.

\bibitem{PlattTrudgian21}
D.J.~Platt and T.~Trudgian,
``The Riemann hypothesis is true up to $3\cdot10^{12}$'',
\emph{Bull.\ Lond.\ Math.\ Soc.} \textbf{53} (2021), no.~3, 792--797.

\bibitem{Connes26}
A.~Connes,
``The Riemann Hypothesis: Past, Present and a Letter Through Time'',
arXiv:2602.04022, 2026.

\bibitem{ConnesVS25}
A.~Connes and W.D.~van Suijlekom,
``Quadratic Forms, Real Zeros and Echoes of the Spectral Action'',
\emph{Comm.\ Math.\ Phys.} \textbf{406} (2025), art.~312,
DOI 10.1007/s00220-025-05493-1; arXiv:2511.23257.

\bibitem{Groskin26a}
A.~Groskin,
``High-Precision Approximation of Riemann Zeros via the Truncated Weil
Form'',
arXiv:2605.20224, 2026.

\bibitem{Montgomery73}
H.L.~Montgomery,
``The pair correlation of zeros of the zeta function'',
in \emph{Analytic Number Theory}, Proc.\ Sympos.\ Pure Math.\ \textbf{24},
Amer.\ Math.\ Soc., 1973, 181--193.

\bibitem{Odlyzko87}
A.M.~Odlyzko,
``On the distribution of spacings between zeros of the zeta function'',
\emph{Math.\ Comp.} \textbf{48} (1987), no.~177, 273--308.

\bibitem{BogomolnyKeating96}
E.B.~Bogomolny and J.P.~Keating,
``Random matrix theory and the Riemann zeros II: $n$-point correlations'',
\emph{Nonlinearity} \textbf{9} (1996), no.~4, 911--935.

\bibitem{RudnickSarnak96}
Z.~Rudnick and P.~Sarnak,
``Zeros of principal $L$-functions and random matrix theory'',
\emph{Duke Math.\ J.} \textbf{81} (1996), no.~2, 269--322.

\bibitem{KatzSarnak99}
N.M.~Katz and P.~Sarnak,
``Zeros of zeta functions and symmetry'',
\emph{Bull.\ Amer.\ Math.\ Soc.} \textbf{36} (1999), no.~1, 1--26.

\bibitem{Besicovitch32}
A.S.~Besicovitch,
\emph{Almost Periodic Functions},
Cambridge University Press, 1932.

\bibitem{Corduneanu68}
C.~Corduneanu,
\emph{Almost Periodic Functions},
Interscience, New York, 1968.

\bibitem{IwaniecKowalski}
H.~Iwaniec and E.~Kowalski,
\emph{Analytic Number Theory},
Amer.\ Math.\ Soc.\ Colloq.\ Publ.\ \textbf{53}, 2004.

\bibitem{Beurling37}
A.~Beurling,
``Analyse de la loi asymptotique de la distribution des nombres premiers
g\'en\'eralis\'es. I'',
\emph{Acta Math.} \textbf{68} (1937), no.~1, 255--291.

\bibitem{DiamondZhang16}
H.G.~Diamond and W.-B.~Zhang,
\emph{Beurling Generalized Numbers},
Math.\ Surveys Monogr.\ \textbf{213}, Amer.\ Math.\ Soc., 2016.

\bibitem{DebruyneVindas17}
G.~Debruyne and J.~Vindas,
``On PNT equivalences for Beurling numbers'',
\emph{Monatsh.\ Math.} \textbf{184} (2017), no.~3, 401--424.

\bibitem{Johansson17}
F.~Johansson,
``Arb: efficient arbitrary-precision midpoint-radius interval arithmetic'',
\emph{IEEE Trans.\ Comput.} \textbf{66} (2017), no.~8, 1281--1292.

\bibitem{mpmath}
The mpmath development team,
\emph{mpmath: a Python library for arbitrary-precision floating-point
arithmetic}, version~1.4.0, 2026, \texttt{http://mpmath.org/}.

\bibitem{DavenportHeilbronn36}
H.~Davenport and H.~Heilbronn,
``On the zeros of certain Dirichlet series'',
\emph{J.\ London Math.\ Soc.} \textbf{11} (1936), 181--185.

\bibitem{PotterTitchmarsh35}
H.S.A.~Potter and E.C.~Titchmarsh,
``The zeros of Epstein's zeta-functions'',
\emph{Proc.\ London Math.\ Soc.} (2) \textbf{39} (1935), 372--384.

\bibitem{Epstein03}
P.~Epstein,
``Zur Theorie allgemeiner Zetafunktionen'',
\emph{Math.\ Ann.} \textbf{56} (1903), 615--644.

\bibitem{Groskin26b}
A.~Groskin,
``A finite Guinand--Weil dictionary and archimedean tail order for the
truncated Weil quadratic form'',
arXiv:2607.02828, 2026.

\bibitem{Spira94}
R.~Spira,
``Some zeros of the Titchmarsh counterexample'',
Math.\ Comp.\ \textbf{63} (1994), no.~208, 747--748.

\bibitem{Suzuki26}
M.~Suzuki,
``Weil's quadratic form via the screw function'',
arXiv:2606.09096, 2026.

\bibitem{Liu26}
K.~Liu,
``A local Hilbert--P\'olya realisation for elliptic curve
$L$-func\-tions'', arXiv:2605.\allowbreak17645v2, 2026.

\bibitem{Kouroupis23}
A.~Kouroupis,
``Composition operators and generalized primes'',
\emph{Proc.\ Amer.\ Math.\ Soc.} \textbf{151} (2023), no.~10, 4291--4305,
DOI 10.1090/proc/16395; arXiv:2208.10170.

\bibitem{HilberdinkPushnitski24}
T.~Hilberdink and A.~Pushnitski,
``Spectral asymptotics for a family of arithmetical matrices and connection
to Beurling primes'',
arXiv:2401.06892, 2024.

\bibitem{Code}
U.~Tehrani,
Verification scripts for the prime--zero coupling series,
GitHub, \url{https://github.com/utehrani/analysislab-nt}.

\end{thebibliography}

% ============================================================
% Footer
% ============================================================
\enlargethispage{3\baselineskip}
\vspace*{\fill}
\begin{minipage}{\linewidth}
\rule{\linewidth}{0.4pt}\smallskip\\
\small\emph{Paper~9 v0.30 \textperiodcentered\ August~2026
\textperiodcentered\ Pauca sed matura}
\end{minipage}

\end{document}
